% Môn: Vật lý
% Lớp: 10
% Chương: Chương III. Động lực học
% Bài: L10C3 Bài 17. Trọng lực và lực căng
% Số câu: 162
% App đọc file này trực tiếp — sửa rồi Commit trên GitHub, bấm Đồng bộ GitHub .tex

% L10C3 Bài 17. Trọng lực và lực căng

% ===== Câu 1 =====
% ID: AIQ_L10C3_FULL_0649
% Mức: NB
\dangbt{Nhận biết trọng lực, trọng lượng và lực căng dây}
\begin{ex}
% ID: AIQ_L10C3_FULL_0649
% Mức: NB
Một vật khối lượng $2\,\text{kg}$ ở gần mặt đất, lấy g= $10\,\text{m}$ /s^2. Tính trọng lượng của vật.
\choice
{\True 20 $N$}
{21 $N$}
{19 $N$}
{40 $N$}
\loigiai{
Chọn A. Trọng lượng P=mg=2·10= $20\,\text{N}$.
}
\end{ex}

% ===== Câu 2 =====
% ID: AIQ_L10C3_FULL_0650
% Mức: NB
\begin{ex}
% ID: AIQ_L10C3_FULL_0650
% Mức: NB
Một vật khối lượng $3\,\text{kg}$ ở gần mặt đất, lấy g= $10\,\text{m}$ /s^2. Tính trọng lượng của vật.
\choice
{30 $N$}
{\True 30 $N$}
{29 $N$}
{60 $N$}
\loigiai{
Chọn B. Trọng lượng P=mg=3·10= $30\,\text{N}$.
}
\end{ex}

% ===== Câu 3 =====
% ID: AIQ_L10C3_FULL_0651
% Mức: NB
\begin{ex}
% ID: AIQ_L10C3_FULL_0651
% Mức: NB
Một vật khối lượng $4\,\text{kg}$ ở gần mặt đất, lấy g= $10\,\text{m}$ /s^2. Tính trọng lượng của vật.
\choice
{40 $N$}
{41 $N$}
{\True 40 $N$}
{80 $N$}
\loigiai{
Chọn C. Trọng lượng P=mg=4·10= $40\,\text{N}$.
}
\end{ex}

% ===== Câu 4 =====
% ID: AIQ_L10C3_FULL_0652
% Mức: NB
\begin{ex}
% ID: AIQ_L10C3_FULL_0652
% Mức: NB
Một vật khối lượng $5\,\text{kg}$ ở gần mặt đất, lấy g= $10\,\text{m}$ /s^2. Tính trọng lượng của vật.
\choice
{50 $N$}
{51 $N$}
{49 $N$}
{\True 50 $N$}
\loigiai{
Chọn D. Trọng lượng P=mg=5·10= $50\,\text{N}$.
}
\end{ex}

% ===== Câu 5 =====
% ID: AIQ_L10C3_FULL_0653
% Mức: TH
\begin{ex}
% ID: AIQ_L10C3_FULL_0653
% Mức: TH
Một vật khối lượng $6\,\text{kg}$ ở gần mặt đất, lấy g= $10\,\text{m}$ /s^2. Tính trọng lượng của vật.
\choice
{\True 60 $N$}
{61 $N$}
{59 $N$}
{120 $N$}
\loigiai{
Chọn A. Trọng lượng P=mg=6·10= $60\,\text{N}$.
}
\end{ex}

% ===== Câu 6 =====
% ID: AIQ_L10C3_FULL_0654
% Mức: TH
\begin{ex}
% ID: AIQ_L10C3_FULL_0654
% Mức: TH
Một vật khối lượng $2\,\text{kg}$ ở gần mặt đất, lấy g= $10\,\text{m}$ /s^2. Tính trọng lượng của vật.
\choice
{20 $N$}
{\True 20 $N$}
{19 $N$}
{40 $N$}
\loigiai{
Chọn B. Trọng lượng P=mg=2·10= $20\,\text{N}$.
}
\end{ex}

% ===== Câu 7 =====
% ID: AIQ_L10C3_FULL_0655
% Mức: TH
\begin{ex}
% ID: AIQ_L10C3_FULL_0655
% Mức: TH
Một vật khối lượng $3\,\text{kg}$ ở gần mặt đất, lấy g= $10\,\text{m}$ /s^2. Tính trọng lượng của vật.
\choice
{30 $N$}
{31 $N$}
{\True 30 $N$}
{60 $N$}
\loigiai{
Chọn C. Trọng lượng P=mg=3·10= $30\,\text{N}$.
}
\end{ex}

% ===== Câu 8 =====
% ID: AIQ_L10C3_FULL_0656
% Mức: TH
\begin{ex}
% ID: AIQ_L10C3_FULL_0656
% Mức: TH
Một vật khối lượng $4\,\text{kg}$ ở gần mặt đất, lấy g= $10\,\text{m}$ /s^2. Tính trọng lượng của vật.
\choice
{40 $N$}
{41 $N$}
{39 $N$}
{\True 40 $N$}
\loigiai{
Chọn D. Trọng lượng P=mg=4·10= $40\,\text{N}$.
}
\end{ex}

% ===== Câu 9 =====
% ID: AIQ_L10C3_FULL_0657
% Mức: VD
\begin{ex}
% ID: AIQ_L10C3_FULL_0657
% Mức: VD
Một vật khối lượng $5\,\text{kg}$ ở gần mặt đất, lấy g= $10\,\text{m}$ /s^2. Tính trọng lượng của vật.
\choice
{\True 50 $N$}
{51 $N$}
{49 $N$}
{100 $N$}
\loigiai{
Chọn A. Trọng lượng P=mg=5·10= $50\,\text{N}$.
}
\end{ex}

% ===== Câu 10 =====
% ID: AIQ_L10C3_FULL_0658
% Mức: VD
\begin{ex}
% ID: AIQ_L10C3_FULL_0658
% Mức: VD
Một vật khối lượng $6\,\text{kg}$ ở gần mặt đất, lấy g= $10\,\text{m}$ /s^2. Tính trọng lượng của vật.
\choice
{60 $N$}
{\True 60 $N$}
{59 $N$}
{120 $N$}
\loigiai{
Chọn B. Trọng lượng P=mg=6·10= $60\,\text{N}$.
}
\end{ex}

% ===== Câu 11 =====
% ID: AIQ_L10C3_FULL_0659
% Mức: TH
\begin{ex}
% ID: AIQ_L10C3_FULL_0659
% Mức: TH
Một vật khối lượng $5\,\text{kg}$ ở gần mặt đất, lấy g= $10\,\text{m}$ /s^2. Tính trọng lượng của vật.
\choiceTF
{\True Giá trị cần tìm là $50\,\text{N}$.}
{Giá trị cần tìm là $51\,\text{N}$.}
{\True Phải xác định đúng các lực tác dụng và chọn chiều dương trước khi lập phương trình.}
{Có thể bỏ qua điều kiện áp dụng và chiều của lực mà kết quả vẫn luôn đúng.}
\loigiai{
\begin{itemchoice}
\item (Đúng) Giá trị cần tìm là $50\,\text{N}$. (Vì): Trọng lượng P=mg=5·10= $50\,\text{N}$. b) Sai vì sai giá trị. c) Đúng. d) Sai vì phải kiểm tra điều kiện áp dụng và chiều lực
\item (Sai) Giá trị cần tìm là $51\,\text{N}$.
\item (Đúng) Phải xác định đúng các lực tác dụng và chọn chiều dương trước khi lập phương trình.
\item (Sai) Có thể bỏ qua điều kiện áp dụng và chiều của lực mà kết quả vẫn luôn đúng.
\end{itemchoice}
}
\end{ex}

% ===== Câu 12 =====
% ID: AIQ_L10C3_FULL_0660
% Mức: TH
\begin{ex}
% ID: AIQ_L10C3_FULL_0660
% Mức: TH
Một vật khối lượng $6\,\text{kg}$ ở gần mặt đất, lấy g= $10\,\text{m}$ /s^2. Tính trọng lượng của vật.
\choiceTF
{\True Giá trị cần tìm là $60\,\text{N}$.}
{Giá trị cần tìm là $61\,\text{N}$.}
{\True Phải xác định đúng các lực tác dụng và chọn chiều dương trước khi lập phương trình.}
{Có thể bỏ qua điều kiện áp dụng và chiều của lực mà kết quả vẫn luôn đúng.}
\loigiai{
\begin{itemchoice}
\item (Đúng) Giá trị cần tìm là $60\,\text{N}$. (Vì): Trọng lượng P=mg=6·10= $60\,\text{N}$. b) Sai vì sai giá trị. c) Đúng. d) Sai vì phải kiểm tra điều kiện áp dụng và chiều lực
\item (Sai) Giá trị cần tìm là $61\,\text{N}$.
\item (Đúng) Phải xác định đúng các lực tác dụng và chọn chiều dương trước khi lập phương trình.
\item (Sai) Có thể bỏ qua điều kiện áp dụng và chiều của lực mà kết quả vẫn luôn đúng.
\end{itemchoice}
}
\end{ex}

% ===== Câu 13 =====
% ID: AIQ_L10C3_FULL_0661
% Mức: VD
\begin{ex}
% ID: AIQ_L10C3_FULL_0661
% Mức: VD
Một vật khối lượng $2\,\text{kg}$ ở gần mặt đất, lấy g= $10\,\text{m}$ /s^2. Tính trọng lượng của vật.
\choiceTF
{\True Giá trị cần tìm là $20\,\text{N}$.}
{Giá trị cần tìm là $21\,\text{N}$.}
{\True Phải xác định đúng các lực tác dụng và chọn chiều dương trước khi lập phương trình.}
{Có thể bỏ qua điều kiện áp dụng và chiều của lực mà kết quả vẫn luôn đúng.}
\loigiai{
\begin{itemchoice}
\item (Đúng) Giá trị cần tìm là $20\,\text{N}$. (Vì): Trọng lượng P=mg=2·10= $20\,\text{N}$. b) Sai vì sai giá trị. c) Đúng. d) Sai vì phải kiểm tra điều kiện áp dụng và chiều lực
\item (Sai) Giá trị cần tìm là $21\,\text{N}$.
\item (Đúng) Phải xác định đúng các lực tác dụng và chọn chiều dương trước khi lập phương trình.
\item (Sai) Có thể bỏ qua điều kiện áp dụng và chiều của lực mà kết quả vẫn luôn đúng.
\end{itemchoice}
}
\end{ex}

% ===== Câu 14 =====
% ID: AIQ_L10C3_FULL_0662
% Mức: VD
\begin{ex}
% ID: AIQ_L10C3_FULL_0662
% Mức: VD
Một vật khối lượng $3\,\text{kg}$ ở gần mặt đất, lấy g= $10\,\text{m}$ /s^2. Tính trọng lượng của vật.
\choiceTF
{\True Giá trị cần tìm là $30\,\text{N}$.}
{Giá trị cần tìm là $31\,\text{N}$.}
{\True Phải xác định đúng các lực tác dụng và chọn chiều dương trước khi lập phương trình.}
{Có thể bỏ qua điều kiện áp dụng và chiều của lực mà kết quả vẫn luôn đúng.}
\loigiai{
\begin{itemchoice}
\item (Đúng) Giá trị cần tìm là $30\,\text{N}$. (Vì): Trọng lượng P=mg=3·10= $30\,\text{N}$. b) Sai vì sai giá trị. c) Đúng. d) Sai vì phải kiểm tra điều kiện áp dụng và chiều lực
\item (Sai) Giá trị cần tìm là $31\,\text{N}$.
\item (Đúng) Phải xác định đúng các lực tác dụng và chọn chiều dương trước khi lập phương trình.
\item (Sai) Có thể bỏ qua điều kiện áp dụng và chiều của lực mà kết quả vẫn luôn đúng.
\end{itemchoice}
}
\end{ex}

% ===== Câu 15 =====
% ID: AIQ_L10C3_FULL_0663
% Mức: VD
\begin{ex}
% ID: AIQ_L10C3_FULL_0663
% Mức: VD
Một vật khối lượng $4\,\text{kg}$ ở gần mặt đất, lấy g= $10\,\text{m}$ /s^2. Tính trọng lượng của vật.
\choiceTF
{\True Giá trị cần tìm là $40\,\text{N}$.}
{Giá trị cần tìm là $41\,\text{N}$.}
{\True Phải xác định đúng các lực tác dụng và chọn chiều dương trước khi lập phương trình.}
{Có thể bỏ qua điều kiện áp dụng và chiều của lực mà kết quả vẫn luôn đúng.}
\loigiai{
\begin{itemchoice}
\item (Đúng) Giá trị cần tìm là $40\,\text{N}$. (Vì): Trọng lượng P=mg=4·10= $40\,\text{N}$. b) Sai vì sai giá trị. c) Đúng. d) Sai vì phải kiểm tra điều kiện áp dụng và chiều lực
\item (Sai) Giá trị cần tìm là $41\,\text{N}$.
\item (Đúng) Phải xác định đúng các lực tác dụng và chọn chiều dương trước khi lập phương trình.
\item (Sai) Có thể bỏ qua điều kiện áp dụng và chiều của lực mà kết quả vẫn luôn đúng.
\end{itemchoice}
}
\end{ex}

% ===== Câu 16 =====
% ID: AIQ_L10C3_FULL_0664
% Mức: TH
\begin{ex}
% ID: AIQ_L10C3_FULL_0664
% Mức: TH
Một vật khối lượng $3\,\text{kg}$ ở gần mặt đất, lấy g= $10\,\text{m}$ /s^2. Tính trọng lượng của vật.
\shortans{30}
\loigiai{
Trọng lượng P=mg=3·10= $30\,\text{N}$.
}
\end{ex}

% ===== Câu 17 =====
% ID: AIQ_L10C3_FULL_0665
% Mức: TH
\begin{ex}
% ID: AIQ_L10C3_FULL_0665
% Mức: TH
Một vật khối lượng $4\,\text{kg}$ ở gần mặt đất, lấy g= $10\,\text{m}$ /s^2. Tính trọng lượng của vật.
\shortans{40}
\loigiai{
Trọng lượng P=mg=4·10= $40\,\text{N}$.
}
\end{ex}

% ===== Câu 18 =====
% ID: AIQ_L10C3_FULL_0666
% Mức: VD
\begin{ex}
% ID: AIQ_L10C3_FULL_0666
% Mức: VD
Một vật khối lượng $5\,\text{kg}$ ở gần mặt đất, lấy g= $10\,\text{m}$ /s^2. Tính trọng lượng của vật.
\shortans{50}
\loigiai{
Trọng lượng P=mg=5·10= $50\,\text{N}$.
}
\end{ex}

% ===== Câu 19 =====
% ID: AIQ_L10C3_FULL_0667
% Mức: VD
\begin{ex}
% ID: AIQ_L10C3_FULL_0667
% Mức: VD
Một vật khối lượng $6\,\text{kg}$ ở gần mặt đất, lấy g= $10\,\text{m}$ /s^2. Tính trọng lượng của vật.
\shortans{60}
\loigiai{
Trọng lượng P=mg=6·10= $60\,\text{N}$.
}
\end{ex}

% ===== Câu 20 =====
% ID: AIQ_L10C3_FULL_0668
% Mức: VD
\begin{ex}
% ID: AIQ_L10C3_FULL_0668
% Mức: VD
Một vật khối lượng $2\,\text{kg}$ ở gần mặt đất, lấy g= $10\,\text{m}$ /s^2. Tính trọng lượng của vật.
\shortans{20}
\loigiai{
Trọng lượng P=mg=2·10= $20\,\text{N}$.
}
\end{ex}

% ===== Câu 21 =====
% ID: AIQ_L10C3_FULL_0669
% Mức: VDC
\begin{ex}
% ID: AIQ_L10C3_FULL_0669
% Mức: VDC
Một vật khối lượng $3\,\text{kg}$ ở gần mặt đất, lấy g= $10\,\text{m}$ /s^2. Tính trọng lượng của vật.
\shortans{30}
\loigiai{
Trọng lượng P=mg=3·10= $30\,\text{N}$.
}
\end{ex}

% ===== Câu 22 =====
% ID: AIQ_L10C3_FULL_0670
% Mức: VD
\begin{bt}
% ID: AIQ_L10C3_FULL_0670
% Mức: VD
Một vật khối lượng $6\,\text{kg}$ ở gần mặt đất, lấy g= $10\,\text{m}$ /s^2. Tính trọng lượng của vật.
\loigiai{
Phân tích lực, chọn hệ quy chiếu và áp dụng định luật hoặc quy tắc phù hợp. Trọng lượng P=mg=6·10= $60\,\text{N}$.
}
\end{bt}

% ===== Câu 23 =====
% ID: AIQ_L10C3_FULL_0671
% Mức: VD
\begin{bt}
% ID: AIQ_L10C3_FULL_0671
% Mức: VD
Một vật khối lượng $2\,\text{kg}$ ở gần mặt đất, lấy g= $10\,\text{m}$ /s^2. Tính trọng lượng của vật.
\loigiai{
Phân tích lực, chọn hệ quy chiếu và áp dụng định luật hoặc quy tắc phù hợp. Trọng lượng P=mg=2·10= $20\,\text{N}$.
}
\end{bt}

% ===== Câu 24 =====
% ID: AIQ_L10C3_FULL_0672
% Mức: VDC
\begin{bt}
% ID: AIQ_L10C3_FULL_0672
% Mức: VDC
Một vật khối lượng $3\,\text{kg}$ ở gần mặt đất, lấy g= $10\,\text{m}$ /s^2. Tính trọng lượng của vật.
\loigiai{
Phân tích lực, chọn hệ quy chiếu và áp dụng định luật hoặc quy tắc phù hợp. Trọng lượng P=mg=3·10= $30\,\text{N}$.
}
\end{bt}

% ===== Câu 25 =====
% ID: AIQ_L10C3_FULL_0673
% Mức: VDC
\begin{bt}
% ID: AIQ_L10C3_FULL_0673
% Mức: VDC
Một vật khối lượng $4\,\text{kg}$ ở gần mặt đất, lấy g= $10\,\text{m}$ /s^2. Tính trọng lượng của vật.
\loigiai{
Phân tích lực, chọn hệ quy chiếu và áp dụng định luật hoặc quy tắc phù hợp. Trọng lượng P=mg=4·10= $40\,\text{N}$.
}
\end{bt}

% ===== Câu 26 =====
% ID: AIQ_L10C3_FULL_0674
% Mức: VDC
\begin{bt}
% ID: AIQ_L10C3_FULL_0674
% Mức: VDC
Một vật khối lượng $5\,\text{kg}$ ở gần mặt đất, lấy g= $10\,\text{m}$ /s^2. Tính trọng lượng của vật.
\loigiai{
Phân tích lực, chọn hệ quy chiếu và áp dụng định luật hoặc quy tắc phù hợp. Trọng lượng P=mg=5·10= $50\,\text{N}$.
}
\end{bt}

% ===== Câu 27 =====
% ID: AIQ_L10C3_FULL_0675
% Mức: VDC
\begin{bt}
% ID: AIQ_L10C3_FULL_0675
% Mức: VDC
Một vật khối lượng $6\,\text{kg}$ ở gần mặt đất, lấy g= $10\,\text{m}$ /s^2. Tính trọng lượng của vật.
\loigiai{
Phân tích lực, chọn hệ quy chiếu và áp dụng định luật hoặc quy tắc phù hợp. Trọng lượng P=mg=6·10= $60\,\text{N}$.
}
\end{bt}

% ===== Câu 28 =====
% ID: AIQ_L10C3_FULL_0676
% Mức: NB
\dangbt{Tính trọng lực và gia tốc rơi tự do}
\begin{ex}
% ID: AIQ_L10C3_FULL_0676
% Mức: NB
Một vật có trọng lượng 20 $N$ tại nơi có gia tốc rơi tự do g= $10\,\text{m}$ /s^2. Tính khối lượng vật.
\choice
{$2\,\text{kg}$}
{\True $2\,\text{kg}$}
{$1\,\text{kg}$}
{$4\,\text{kg}$}
\loigiai{
Chọn B. m=P/g=20/10= $2\,\text{kg}$.
}
\end{ex}

% ===== Câu 29 =====
% ID: AIQ_L10C3_FULL_0677
% Mức: NB
\begin{ex}
% ID: AIQ_L10C3_FULL_0677
% Mức: NB
Một vật có trọng lượng 30 $N$ tại nơi có gia tốc rơi tự do g= $10\,\text{m}$ /s^2. Tính khối lượng vật.
\choice
{$3\,\text{kg}$}
{$4\,\text{kg}$}
{\True $3\,\text{kg}$}
{$6\,\text{kg}$}
\loigiai{
Chọn C. m=P/g=30/10= $3\,\text{kg}$.
}
\end{ex}

% ===== Câu 30 =====
% ID: AIQ_L10C3_FULL_0678
% Mức: NB
\begin{ex}
% ID: AIQ_L10C3_FULL_0678
% Mức: NB
Một vật có trọng lượng 40 $N$ tại nơi có gia tốc rơi tự do g= $10\,\text{m}$ /s^2. Tính khối lượng vật.
\choice
{$4\,\text{kg}$}
{$5\,\text{kg}$}
{$3\,\text{kg}$}
{\True $4\,\text{kg}$}
\loigiai{
Chọn D. m=P/g=40/10= $4\,\text{kg}$.
}
\end{ex}

% ===== Câu 31 =====
% ID: AIQ_L10C3_FULL_0679
% Mức: NB
\begin{ex}
% ID: AIQ_L10C3_FULL_0679
% Mức: NB
Một vật có trọng lượng 50 $N$ tại nơi có gia tốc rơi tự do g= $10\,\text{m}$ /s^2. Tính khối lượng vật.
\choice
{\True $5\,\text{kg}$}
{$6\,\text{kg}$}
{$4\,\text{kg}$}
{$10\,\text{kg}$}
\loigiai{
Chọn A. m=P/g=50/10= $5\,\text{kg}$.
}
\end{ex}

% ===== Câu 32 =====
% ID: AIQ_L10C3_FULL_0680
% Mức: TH
\begin{ex}
% ID: AIQ_L10C3_FULL_0680
% Mức: TH
Một vật có trọng lượng 60 $N$ tại nơi có gia tốc rơi tự do g= $10\,\text{m}$ /s^2. Tính khối lượng vật.
\choice
{$6\,\text{kg}$}
{\True $6\,\text{kg}$}
{$5\,\text{kg}$}
{$12\,\text{kg}$}
\loigiai{
Chọn B. m=P/g=60/10= $6\,\text{kg}$.
}
\end{ex}

% ===== Câu 33 =====
% ID: AIQ_L10C3_FULL_0681
% Mức: TH
\begin{ex}
% ID: AIQ_L10C3_FULL_0681
% Mức: TH
Một vật có trọng lượng 20 $N$ tại nơi có gia tốc rơi tự do g= $10\,\text{m}$ /s^2. Tính khối lượng vật.
\choice
{$2\,\text{kg}$}
{$3\,\text{kg}$}
{\True $2\,\text{kg}$}
{$4\,\text{kg}$}
\loigiai{
Chọn C. m=P/g=20/10= $2\,\text{kg}$.
}
\end{ex}

% ===== Câu 34 =====
% ID: AIQ_L10C3_FULL_0682
% Mức: TH
\begin{ex}
% ID: AIQ_L10C3_FULL_0682
% Mức: TH
Một vật có trọng lượng 30 $N$ tại nơi có gia tốc rơi tự do g= $10\,\text{m}$ /s^2. Tính khối lượng vật.
\choice
{$3\,\text{kg}$}
{$4\,\text{kg}$}
{$2\,\text{kg}$}
{\True $3\,\text{kg}$}
\loigiai{
Chọn D. m=P/g=30/10= $3\,\text{kg}$.
}
\end{ex}

% ===== Câu 35 =====
% ID: AIQ_L10C3_FULL_0683
% Mức: TH
\begin{ex}
% ID: AIQ_L10C3_FULL_0683
% Mức: TH
Một vật có trọng lượng 40 $N$ tại nơi có gia tốc rơi tự do g= $10\,\text{m}$ /s^2. Tính khối lượng vật.
\choice
{\True $4\,\text{kg}$}
{$5\,\text{kg}$}
{$3\,\text{kg}$}
{$8\,\text{kg}$}
\loigiai{
Chọn A. m=P/g=40/10= $4\,\text{kg}$.
}
\end{ex}

% ===== Câu 36 =====
% ID: AIQ_L10C3_FULL_0684
% Mức: VD
\begin{ex}
% ID: AIQ_L10C3_FULL_0684
% Mức: VD
Một vật có trọng lượng 50 $N$ tại nơi có gia tốc rơi tự do g= $10\,\text{m}$ /s^2. Tính khối lượng vật.
\choice
{$5\,\text{kg}$}
{\True $5\,\text{kg}$}
{$4\,\text{kg}$}
{$10\,\text{kg}$}
\loigiai{
Chọn B. m=P/g=50/10= $5\,\text{kg}$.
}
\end{ex}

% ===== Câu 37 =====
% ID: AIQ_L10C3_FULL_0685
% Mức: VD
\begin{ex}
% ID: AIQ_L10C3_FULL_0685
% Mức: VD
Một vật có trọng lượng 60 $N$ tại nơi có gia tốc rơi tự do g= $10\,\text{m}$ /s^2. Tính khối lượng vật.
\choice
{$6\,\text{kg}$}
{$7\,\text{kg}$}
{\True $6\,\text{kg}$}
{$12\,\text{kg}$}
\loigiai{
Chọn C. m=P/g=60/10= $6\,\text{kg}$.
}
\end{ex}

% ===== Câu 38 =====
% ID: AIQ_L10C3_FULL_0686
% Mức: TH
\begin{ex}
% ID: AIQ_L10C3_FULL_0686
% Mức: TH
Một vật có trọng lượng 50 $N$ tại nơi có gia tốc rơi tự do g= $10\,\text{m}$ /s^2. Tính khối lượng vật.
\choiceTF
{\True Giá trị cần tìm là $5\,\text{kg}$.}
{Giá trị cần tìm là $6\,\text{kg}$.}
{\True Phải xác định đúng các lực tác dụng và chọn chiều dương trước khi lập phương trình.}
{Có thể bỏ qua điều kiện áp dụng và chiều của lực mà kết quả vẫn luôn đúng.}
\loigiai{
\begin{itemchoice}
\item (Đúng) Giá trị cần tìm là $5\,\text{kg}$. (Vì): m=P/g=50/10= $5\,\text{kg}$. b) Sai vì sai giá trị. c) Đúng. d) Sai vì phải kiểm tra điều kiện áp dụng và chiều lực
\item (Sai) Giá trị cần tìm là $6\,\text{kg}$.
\item (Đúng) Phải xác định đúng các lực tác dụng và chọn chiều dương trước khi lập phương trình.
\item (Sai) Có thể bỏ qua điều kiện áp dụng và chiều của lực mà kết quả vẫn luôn đúng.
\end{itemchoice}
}
\end{ex}

% ===== Câu 39 =====
% ID: AIQ_L10C3_FULL_0687
% Mức: TH
\begin{ex}
% ID: AIQ_L10C3_FULL_0687
% Mức: TH
Một vật có trọng lượng 60 $N$ tại nơi có gia tốc rơi tự do g= $10\,\text{m}$ /s^2. Tính khối lượng vật.
\choiceTF
{\True Giá trị cần tìm là $6\,\text{kg}$.}
{Giá trị cần tìm là $7\,\text{kg}$.}
{\True Phải xác định đúng các lực tác dụng và chọn chiều dương trước khi lập phương trình.}
{Có thể bỏ qua điều kiện áp dụng và chiều của lực mà kết quả vẫn luôn đúng.}
\loigiai{
\begin{itemchoice}
\item (Đúng) Giá trị cần tìm là $6\,\text{kg}$. (Vì): m=P/g=60/10= $6\,\text{kg}$. b) Sai vì sai giá trị. c) Đúng. d) Sai vì phải kiểm tra điều kiện áp dụng và chiều lực
\item (Sai) Giá trị cần tìm là $7\,\text{kg}$.
\item (Đúng) Phải xác định đúng các lực tác dụng và chọn chiều dương trước khi lập phương trình.
\item (Sai) Có thể bỏ qua điều kiện áp dụng và chiều của lực mà kết quả vẫn luôn đúng.
\end{itemchoice}
}
\end{ex}

% ===== Câu 40 =====
% ID: AIQ_L10C3_FULL_0688
% Mức: VD
\begin{ex}
% ID: AIQ_L10C3_FULL_0688
% Mức: VD
Một vật có trọng lượng 20 $N$ tại nơi có gia tốc rơi tự do g= $10\,\text{m}$ /s^2. Tính khối lượng vật.
\choiceTF
{\True Giá trị cần tìm là $2\,\text{kg}$.}
{Giá trị cần tìm là $3\,\text{kg}$.}
{\True Phải xác định đúng các lực tác dụng và chọn chiều dương trước khi lập phương trình.}
{Có thể bỏ qua điều kiện áp dụng và chiều của lực mà kết quả vẫn luôn đúng.}
\loigiai{
\begin{itemchoice}
\item (Đúng) Giá trị cần tìm là $2\,\text{kg}$. (Vì): m=P/g=20/10= $2\,\text{kg}$. b) Sai vì sai giá trị. c) Đúng. d) Sai vì phải kiểm tra điều kiện áp dụng và chiều lực
\item (Sai) Giá trị cần tìm là $3\,\text{kg}$.
\item (Đúng) Phải xác định đúng các lực tác dụng và chọn chiều dương trước khi lập phương trình.
\item (Sai) Có thể bỏ qua điều kiện áp dụng và chiều của lực mà kết quả vẫn luôn đúng.
\end{itemchoice}
}
\end{ex}

% ===== Câu 41 =====
% ID: AIQ_L10C3_FULL_0689
% Mức: VD
\begin{ex}
% ID: AIQ_L10C3_FULL_0689
% Mức: VD
Một vật có trọng lượng 30 $N$ tại nơi có gia tốc rơi tự do g= $10\,\text{m}$ /s^2. Tính khối lượng vật.
\choiceTF
{\True Giá trị cần tìm là $3\,\text{kg}$.}
{Giá trị cần tìm là $4\,\text{kg}$.}
{\True Phải xác định đúng các lực tác dụng và chọn chiều dương trước khi lập phương trình.}
{Có thể bỏ qua điều kiện áp dụng và chiều của lực mà kết quả vẫn luôn đúng.}
\loigiai{
\begin{itemchoice}
\item (Đúng) Giá trị cần tìm là $3\,\text{kg}$. (Vì): m=P/g=30/10= $3\,\text{kg}$. b) Sai vì sai giá trị. c) Đúng. d) Sai vì phải kiểm tra điều kiện áp dụng và chiều lực
\item (Sai) Giá trị cần tìm là $4\,\text{kg}$.
\item (Đúng) Phải xác định đúng các lực tác dụng và chọn chiều dương trước khi lập phương trình.
\item (Sai) Có thể bỏ qua điều kiện áp dụng và chiều của lực mà kết quả vẫn luôn đúng.
\end{itemchoice}
}
\end{ex}

% ===== Câu 42 =====
% ID: AIQ_L10C3_FULL_0690
% Mức: VD
\begin{ex}
% ID: AIQ_L10C3_FULL_0690
% Mức: VD
Một vật có trọng lượng 40 $N$ tại nơi có gia tốc rơi tự do g= $10\,\text{m}$ /s^2. Tính khối lượng vật.
\choiceTF
{\True Giá trị cần tìm là $4\,\text{kg}$.}
{Giá trị cần tìm là $5\,\text{kg}$.}
{\True Phải xác định đúng các lực tác dụng và chọn chiều dương trước khi lập phương trình.}
{Có thể bỏ qua điều kiện áp dụng và chiều của lực mà kết quả vẫn luôn đúng.}
\loigiai{
\begin{itemchoice}
\item (Đúng) Giá trị cần tìm là $4\,\text{kg}$. (Vì): m=P/g=40/10= $4\,\text{kg}$. b) Sai vì sai giá trị. c) Đúng. d) Sai vì phải kiểm tra điều kiện áp dụng và chiều lực
\item (Sai) Giá trị cần tìm là $5\,\text{kg}$.
\item (Đúng) Phải xác định đúng các lực tác dụng và chọn chiều dương trước khi lập phương trình.
\item (Sai) Có thể bỏ qua điều kiện áp dụng và chiều của lực mà kết quả vẫn luôn đúng.
\end{itemchoice}
}
\end{ex}

% ===== Câu 43 =====
% ID: AIQ_L10C3_FULL_0691
% Mức: TH
\begin{ex}
% ID: AIQ_L10C3_FULL_0691
% Mức: TH
Một vật có trọng lượng 30 $N$ tại nơi có gia tốc rơi tự do g= $10\,\text{m}$ /s^2. Tính khối lượng vật.
\shortans{3}
\loigiai{
$m=P/g=30/10=3 kg.$
}
\end{ex}

% ===== Câu 44 =====
% ID: AIQ_L10C3_FULL_0692
% Mức: TH
\begin{ex}
% ID: AIQ_L10C3_FULL_0692
% Mức: TH
Một vật có trọng lượng 40 $N$ tại nơi có gia tốc rơi tự do g= $10\,\text{m}$ /s^2. Tính khối lượng vật.
\shortans{4}
\loigiai{
$m=P/g=40/10=4 kg.$
}
\end{ex}

% ===== Câu 45 =====
% ID: AIQ_L10C3_FULL_0693
% Mức: VD
\begin{ex}
% ID: AIQ_L10C3_FULL_0693
% Mức: VD
Một vật có trọng lượng 50 $N$ tại nơi có gia tốc rơi tự do g= $10\,\text{m}$ /s^2. Tính khối lượng vật.
\shortans{5}
\loigiai{
$m=P/g=50/10=5 kg.$
}
\end{ex}

% ===== Câu 46 =====
% ID: AIQ_L10C3_FULL_0694
% Mức: VD
\begin{ex}
% ID: AIQ_L10C3_FULL_0694
% Mức: VD
Một vật có trọng lượng 60 $N$ tại nơi có gia tốc rơi tự do g= $10\,\text{m}$ /s^2. Tính khối lượng vật.
\shortans{6}
\loigiai{
$m=P/g=60/10=6 kg.$
}
\end{ex}

% ===== Câu 47 =====
% ID: AIQ_L10C3_FULL_0695
% Mức: VD
\begin{ex}
% ID: AIQ_L10C3_FULL_0695
% Mức: VD
Một vật có trọng lượng 20 $N$ tại nơi có gia tốc rơi tự do g= $10\,\text{m}$ /s^2. Tính khối lượng vật.
\shortans{2}
\loigiai{
$m=P/g=20/10=2 kg.$
}
\end{ex}

% ===== Câu 48 =====
% ID: AIQ_L10C3_FULL_0696
% Mức: VDC
\begin{ex}
% ID: AIQ_L10C3_FULL_0696
% Mức: VDC
Một vật có trọng lượng 30 $N$ tại nơi có gia tốc rơi tự do g= $10\,\text{m}$ /s^2. Tính khối lượng vật.
\shortans{3}
\loigiai{
$m=P/g=30/10=3 kg.$
}
\end{ex}

% ===== Câu 49 =====
% ID: AIQ_L10C3_FULL_0697
% Mức: VD
\begin{bt}
% ID: AIQ_L10C3_FULL_0697
% Mức: VD
Một vật có trọng lượng 60 $N$ tại nơi có gia tốc rơi tự do g= $10\,\text{m}$ /s^2. Tính khối lượng vật.
\loigiai{
Phân tích lực, chọn hệ quy chiếu và áp dụng định luật hoặc quy tắc phù hợp. m=P/g=60/10= $6\,\text{kg}$.
}
\end{bt}

% ===== Câu 50 =====
% ID: AIQ_L10C3_FULL_0698
% Mức: VD
\begin{bt}
% ID: AIQ_L10C3_FULL_0698
% Mức: VD
Một vật có trọng lượng 20 $N$ tại nơi có gia tốc rơi tự do g= $10\,\text{m}$ /s^2. Tính khối lượng vật.
\loigiai{
Phân tích lực, chọn hệ quy chiếu và áp dụng định luật hoặc quy tắc phù hợp. m=P/g=20/10= $2\,\text{kg}$.
}
\end{bt}

% ===== Câu 51 =====
% ID: AIQ_L10C3_FULL_0699
% Mức: VDC
\begin{bt}
% ID: AIQ_L10C3_FULL_0699
% Mức: VDC
Một vật có trọng lượng 30 $N$ tại nơi có gia tốc rơi tự do g= $10\,\text{m}$ /s^2. Tính khối lượng vật.
\loigiai{
Phân tích lực, chọn hệ quy chiếu và áp dụng định luật hoặc quy tắc phù hợp. m=P/g=30/10= $3\,\text{kg}$.
}
\end{bt}

% ===== Câu 52 =====
% ID: AIQ_L10C3_FULL_0700
% Mức: VDC
\begin{bt}
% ID: AIQ_L10C3_FULL_0700
% Mức: VDC
Một vật có trọng lượng 40 $N$ tại nơi có gia tốc rơi tự do g= $10\,\text{m}$ /s^2. Tính khối lượng vật.
\loigiai{
Phân tích lực, chọn hệ quy chiếu và áp dụng định luật hoặc quy tắc phù hợp. m=P/g=40/10= $4\,\text{kg}$.
}
\end{bt}

% ===== Câu 53 =====
% ID: AIQ_L10C3_FULL_0701
% Mức: VDC
\begin{bt}
% ID: AIQ_L10C3_FULL_0701
% Mức: VDC
Một vật có trọng lượng 50 $N$ tại nơi có gia tốc rơi tự do g= $10\,\text{m}$ /s^2. Tính khối lượng vật.
\loigiai{
Phân tích lực, chọn hệ quy chiếu và áp dụng định luật hoặc quy tắc phù hợp. m=P/g=50/10= $5\,\text{kg}$.
}
\end{bt}

% ===== Câu 54 =====
% ID: AIQ_L10C3_FULL_0702
% Mức: VDC
\begin{bt}
% ID: AIQ_L10C3_FULL_0702
% Mức: VDC
Một vật có trọng lượng 60 $N$ tại nơi có gia tốc rơi tự do g= $10\,\text{m}$ /s^2. Tính khối lượng vật.
\loigiai{
Phân tích lực, chọn hệ quy chiếu và áp dụng định luật hoặc quy tắc phù hợp. m=P/g=60/10= $6\,\text{kg}$.
}
\end{bt}

% ===== Câu 55 =====
% ID: AIQ_L10C3_FULL_0703
% Mức: NB
\dangbt{Xác định lực căng của dây treo vật cân bằng}
\begin{ex}
% ID: AIQ_L10C3_FULL_0703
% Mức: NB
Vật $2\,\text{kg}$ treo đứng yên bằng một sợi dây nhẹ thẳng đứng. Lấy g= $10\,\text{m}$ /s^2. Tính lực căng dây.
\choice
{20 $N$}
{21 $N$}
{\True 20 $N$}
{40 $N$}
\loigiai{
Chọn C. Vật cân bằng nên T=P=mg= $20\,\text{N}$.
}
\end{ex}

% ===== Câu 56 =====
% ID: AIQ_L10C3_FULL_0704
% Mức: NB
\begin{ex}
% ID: AIQ_L10C3_FULL_0704
% Mức: NB
Vật $3\,\text{kg}$ treo đứng yên bằng một sợi dây nhẹ thẳng đứng. Lấy g= $10\,\text{m}$ /s^2. Tính lực căng dây.
\choice
{30 $N$}
{31 $N$}
{29 $N$}
{\True 30 $N$}
\loigiai{
Chọn D. Vật cân bằng nên T=P=mg= $30\,\text{N}$.
}
\end{ex}

% ===== Câu 57 =====
% ID: AIQ_L10C3_FULL_0705
% Mức: NB
\begin{ex}
% ID: AIQ_L10C3_FULL_0705
% Mức: NB
Vật $4\,\text{kg}$ treo đứng yên bằng một sợi dây nhẹ thẳng đứng. Lấy g= $10\,\text{m}$ /s^2. Tính lực căng dây.
\choice
{\True 40 $N$}
{41 $N$}
{39 $N$}
{80 $N$}
\loigiai{
Chọn A. Vật cân bằng nên T=P=mg= $40\,\text{N}$.
}
\end{ex}

% ===== Câu 58 =====
% ID: AIQ_L10C3_FULL_0706
% Mức: NB
\begin{ex}
% ID: AIQ_L10C3_FULL_0706
% Mức: NB
Vật $5\,\text{kg}$ treo đứng yên bằng một sợi dây nhẹ thẳng đứng. Lấy g= $10\,\text{m}$ /s^2. Tính lực căng dây.
\choice
{50 $N$}
{\True 50 $N$}
{49 $N$}
{100 $N$}
\loigiai{
Chọn B. Vật cân bằng nên T=P=mg= $50\,\text{N}$.
}
\end{ex}

% ===== Câu 59 =====
% ID: AIQ_L10C3_FULL_0707
% Mức: TH
\begin{ex}
% ID: AIQ_L10C3_FULL_0707
% Mức: TH
Vật $6\,\text{kg}$ treo đứng yên bằng một sợi dây nhẹ thẳng đứng. Lấy g= $10\,\text{m}$ /s^2. Tính lực căng dây.
\choice
{60 $N$}
{61 $N$}
{\True 60 $N$}
{120 $N$}
\loigiai{
Chọn C. Vật cân bằng nên T=P=mg= $60\,\text{N}$.
}
\end{ex}

% ===== Câu 60 =====
% ID: AIQ_L10C3_FULL_0708
% Mức: TH
\begin{ex}
% ID: AIQ_L10C3_FULL_0708
% Mức: TH
Vật $2\,\text{kg}$ treo đứng yên bằng một sợi dây nhẹ thẳng đứng. Lấy g= $10\,\text{m}$ /s^2. Tính lực căng dây.
\choice
{20 $N$}
{21 $N$}
{19 $N$}
{\True 20 $N$}
\loigiai{
Chọn D. Vật cân bằng nên T=P=mg= $20\,\text{N}$.
}
\end{ex}

% ===== Câu 61 =====
% ID: AIQ_L10C3_FULL_0709
% Mức: TH
\begin{ex}
% ID: AIQ_L10C3_FULL_0709
% Mức: TH
Vật $3\,\text{kg}$ treo đứng yên bằng một sợi dây nhẹ thẳng đứng. Lấy g= $10\,\text{m}$ /s^2. Tính lực căng dây.
\choice
{\True 30 $N$}
{31 $N$}
{29 $N$}
{60 $N$}
\loigiai{
Chọn A. Vật cân bằng nên T=P=mg= $30\,\text{N}$.
}
\end{ex}

% ===== Câu 62 =====
% ID: AIQ_L10C3_FULL_0710
% Mức: TH
\begin{ex}
% ID: AIQ_L10C3_FULL_0710
% Mức: TH
Vật $4\,\text{kg}$ treo đứng yên bằng một sợi dây nhẹ thẳng đứng. Lấy g= $10\,\text{m}$ /s^2. Tính lực căng dây.
\choice
{40 $N$}
{\True 40 $N$}
{39 $N$}
{80 $N$}
\loigiai{
Chọn B. Vật cân bằng nên T=P=mg= $40\,\text{N}$.
}
\end{ex}

% ===== Câu 63 =====
% ID: AIQ_L10C3_FULL_0711
% Mức: VD
\begin{ex}
% ID: AIQ_L10C3_FULL_0711
% Mức: VD
Vật $5\,\text{kg}$ treo đứng yên bằng một sợi dây nhẹ thẳng đứng. Lấy g= $10\,\text{m}$ /s^2. Tính lực căng dây.
\choice
{50 $N$}
{51 $N$}
{\True 50 $N$}
{100 $N$}
\loigiai{
Chọn C. Vật cân bằng nên T=P=mg= $50\,\text{N}$.
}
\end{ex}

% ===== Câu 64 =====
% ID: AIQ_L10C3_FULL_0712
% Mức: VD
\begin{ex}
% ID: AIQ_L10C3_FULL_0712
% Mức: VD
Vật $6\,\text{kg}$ treo đứng yên bằng một sợi dây nhẹ thẳng đứng. Lấy g= $10\,\text{m}$ /s^2. Tính lực căng dây.
\choice
{60 $N$}
{61 $N$}
{59 $N$}
{\True 60 $N$}
\loigiai{
Chọn D. Vật cân bằng nên T=P=mg= $60\,\text{N}$.
}
\end{ex}

% ===== Câu 65 =====
% ID: AIQ_L10C3_FULL_0713
% Mức: TH
\begin{ex}
% ID: AIQ_L10C3_FULL_0713
% Mức: TH
Vật $5\,\text{kg}$ treo đứng yên bằng một sợi dây nhẹ thẳng đứng. Lấy g= $10\,\text{m}$ /s^2. Tính lực căng dây.
\choiceTF
{\True Giá trị cần tìm là $50\,\text{N}$.}
{Giá trị cần tìm là $51\,\text{N}$.}
{\True Phải xác định đúng các lực tác dụng và chọn chiều dương trước khi lập phương trình.}
{Có thể bỏ qua điều kiện áp dụng và chiều của lực mà kết quả vẫn luôn đúng.}
\loigiai{
\begin{itemchoice}
\item (Đúng) Giá trị cần tìm là $50\,\text{N}$. (Vì): Vật cân bằng nên T=P=mg= $50\,\text{N}$. b) Sai vì sai giá trị. c) Đúng. d) Sai vì phải kiểm tra điều kiện áp dụng và chiều lực
\item (Sai) Giá trị cần tìm là $51\,\text{N}$.
\item (Đúng) Phải xác định đúng các lực tác dụng và chọn chiều dương trước khi lập phương trình.
\item (Sai) Có thể bỏ qua điều kiện áp dụng và chiều của lực mà kết quả vẫn luôn đúng.
\end{itemchoice}
}
\end{ex}

% ===== Câu 66 =====
% ID: AIQ_L10C3_FULL_0714
% Mức: TH
\begin{ex}
% ID: AIQ_L10C3_FULL_0714
% Mức: TH
Vật $6\,\text{kg}$ treo đứng yên bằng một sợi dây nhẹ thẳng đứng. Lấy g= $10\,\text{m}$ /s^2. Tính lực căng dây.
\choiceTF
{\True Giá trị cần tìm là $60\,\text{N}$.}
{Giá trị cần tìm là $61\,\text{N}$.}
{\True Phải xác định đúng các lực tác dụng và chọn chiều dương trước khi lập phương trình.}
{Có thể bỏ qua điều kiện áp dụng và chiều của lực mà kết quả vẫn luôn đúng.}
\loigiai{
\begin{itemchoice}
\item (Đúng) Giá trị cần tìm là $60\,\text{N}$. (Vì): Vật cân bằng nên T=P=mg= $60\,\text{N}$. b) Sai vì sai giá trị. c) Đúng. d) Sai vì phải kiểm tra điều kiện áp dụng và chiều lực
\item (Sai) Giá trị cần tìm là $61\,\text{N}$.
\item (Đúng) Phải xác định đúng các lực tác dụng và chọn chiều dương trước khi lập phương trình.
\item (Sai) Có thể bỏ qua điều kiện áp dụng và chiều của lực mà kết quả vẫn luôn đúng.
\end{itemchoice}
}
\end{ex}

% ===== Câu 67 =====
% ID: AIQ_L10C3_FULL_0715
% Mức: VD
\begin{ex}
% ID: AIQ_L10C3_FULL_0715
% Mức: VD
Vật $2\,\text{kg}$ treo đứng yên bằng một sợi dây nhẹ thẳng đứng. Lấy g= $10\,\text{m}$ /s^2. Tính lực căng dây.
\choiceTF
{\True Giá trị cần tìm là $20\,\text{N}$.}
{Giá trị cần tìm là $21\,\text{N}$.}
{\True Phải xác định đúng các lực tác dụng và chọn chiều dương trước khi lập phương trình.}
{Có thể bỏ qua điều kiện áp dụng và chiều của lực mà kết quả vẫn luôn đúng.}
\loigiai{
\begin{itemchoice}
\item (Đúng) Giá trị cần tìm là $20\,\text{N}$. (Vì): Vật cân bằng nên T=P=mg= $20\,\text{N}$. b) Sai vì sai giá trị. c) Đúng. d) Sai vì phải kiểm tra điều kiện áp dụng và chiều lực
\item (Sai) Giá trị cần tìm là $21\,\text{N}$.
\item (Đúng) Phải xác định đúng các lực tác dụng và chọn chiều dương trước khi lập phương trình.
\item (Sai) Có thể bỏ qua điều kiện áp dụng và chiều của lực mà kết quả vẫn luôn đúng.
\end{itemchoice}
}
\end{ex}

% ===== Câu 68 =====
% ID: AIQ_L10C3_FULL_0716
% Mức: VD
\begin{ex}
% ID: AIQ_L10C3_FULL_0716
% Mức: VD
Vật $3\,\text{kg}$ treo đứng yên bằng một sợi dây nhẹ thẳng đứng. Lấy g= $10\,\text{m}$ /s^2. Tính lực căng dây.
\choiceTF
{\True Giá trị cần tìm là $30\,\text{N}$.}
{Giá trị cần tìm là $31\,\text{N}$.}
{\True Phải xác định đúng các lực tác dụng và chọn chiều dương trước khi lập phương trình.}
{Có thể bỏ qua điều kiện áp dụng và chiều của lực mà kết quả vẫn luôn đúng.}
\loigiai{
\begin{itemchoice}
\item (Đúng) Giá trị cần tìm là $30\,\text{N}$. (Vì): Vật cân bằng nên T=P=mg= $30\,\text{N}$. b) Sai vì sai giá trị. c) Đúng. d) Sai vì phải kiểm tra điều kiện áp dụng và chiều lực
\item (Sai) Giá trị cần tìm là $31\,\text{N}$.
\item (Đúng) Phải xác định đúng các lực tác dụng và chọn chiều dương trước khi lập phương trình.
\item (Sai) Có thể bỏ qua điều kiện áp dụng và chiều của lực mà kết quả vẫn luôn đúng.
\end{itemchoice}
}
\end{ex}

% ===== Câu 69 =====
% ID: AIQ_L10C3_FULL_0717
% Mức: VD
\begin{ex}
% ID: AIQ_L10C3_FULL_0717
% Mức: VD
Vật $4\,\text{kg}$ treo đứng yên bằng một sợi dây nhẹ thẳng đứng. Lấy g= $10\,\text{m}$ /s^2. Tính lực căng dây.
\choiceTF
{\True Giá trị cần tìm là $40\,\text{N}$.}
{Giá trị cần tìm là $41\,\text{N}$.}
{\True Phải xác định đúng các lực tác dụng và chọn chiều dương trước khi lập phương trình.}
{Có thể bỏ qua điều kiện áp dụng và chiều của lực mà kết quả vẫn luôn đúng.}
\loigiai{
\begin{itemchoice}
\item (Đúng) Giá trị cần tìm là $40\,\text{N}$. (Vì): Vật cân bằng nên T=P=mg= $40\,\text{N}$. b) Sai vì sai giá trị. c) Đúng. d) Sai vì phải kiểm tra điều kiện áp dụng và chiều lực
\item (Sai) Giá trị cần tìm là $41\,\text{N}$.
\item (Đúng) Phải xác định đúng các lực tác dụng và chọn chiều dương trước khi lập phương trình.
\item (Sai) Có thể bỏ qua điều kiện áp dụng và chiều của lực mà kết quả vẫn luôn đúng.
\end{itemchoice}
}
\end{ex}

% ===== Câu 70 =====
% ID: AIQ_L10C3_FULL_0718
% Mức: TH
\begin{ex}
% ID: AIQ_L10C3_FULL_0718
% Mức: TH
Vật $3\,\text{kg}$ treo đứng yên bằng một sợi dây nhẹ thẳng đứng. Lấy g= $10\,\text{m}$ /s^2. Tính lực căng dây.
\shortans{30}
\loigiai{
Vật cân bằng nên T=P=mg= $30\,\text{N}$.
}
\end{ex}

% ===== Câu 71 =====
% ID: AIQ_L10C3_FULL_0719
% Mức: TH
\begin{ex}
% ID: AIQ_L10C3_FULL_0719
% Mức: TH
Vật $4\,\text{kg}$ treo đứng yên bằng một sợi dây nhẹ thẳng đứng. Lấy g= $10\,\text{m}$ /s^2. Tính lực căng dây.
\shortans{40}
\loigiai{
Vật cân bằng nên T=P=mg= $40\,\text{N}$.
}
\end{ex}

% ===== Câu 72 =====
% ID: AIQ_L10C3_FULL_0720
% Mức: VD
\begin{ex}
% ID: AIQ_L10C3_FULL_0720
% Mức: VD
Vật $5\,\text{kg}$ treo đứng yên bằng một sợi dây nhẹ thẳng đứng. Lấy g= $10\,\text{m}$ /s^2. Tính lực căng dây.
\shortans{50}
\loigiai{
Vật cân bằng nên T=P=mg= $50\,\text{N}$.
}
\end{ex}

% ===== Câu 73 =====
% ID: AIQ_L10C3_FULL_0721
% Mức: VD
\begin{ex}
% ID: AIQ_L10C3_FULL_0721
% Mức: VD
Vật $6\,\text{kg}$ treo đứng yên bằng một sợi dây nhẹ thẳng đứng. Lấy g= $10\,\text{m}$ /s^2. Tính lực căng dây.
\shortans{60}
\loigiai{
Vật cân bằng nên T=P=mg= $60\,\text{N}$.
}
\end{ex}

% ===== Câu 74 =====
% ID: AIQ_L10C3_FULL_0722
% Mức: VD
\begin{ex}
% ID: AIQ_L10C3_FULL_0722
% Mức: VD
Vật $2\,\text{kg}$ treo đứng yên bằng một sợi dây nhẹ thẳng đứng. Lấy g= $10\,\text{m}$ /s^2. Tính lực căng dây.
\shortans{20}
\loigiai{
Vật cân bằng nên T=P=mg= $20\,\text{N}$.
}
\end{ex}

% ===== Câu 75 =====
% ID: AIQ_L10C3_FULL_0723
% Mức: VDC
\begin{ex}
% ID: AIQ_L10C3_FULL_0723
% Mức: VDC
Vật $3\,\text{kg}$ treo đứng yên bằng một sợi dây nhẹ thẳng đứng. Lấy g= $10\,\text{m}$ /s^2. Tính lực căng dây.
\shortans{30}
\loigiai{
Vật cân bằng nên T=P=mg= $30\,\text{N}$.
}
\end{ex}

% ===== Câu 76 =====
% ID: AIQ_L10C3_FULL_0724
% Mức: VD
\begin{bt}
% ID: AIQ_L10C3_FULL_0724
% Mức: VD
Vật $6\,\text{kg}$ treo đứng yên bằng một sợi dây nhẹ thẳng đứng. Lấy g= $10\,\text{m}$ /s^2. Tính lực căng dây.
\loigiai{
Phân tích lực, chọn hệ quy chiếu và áp dụng định luật hoặc quy tắc phù hợp. Vật cân bằng nên T=P=mg= $60\,\text{N}$.
}
\end{bt}

% ===== Câu 77 =====
% ID: AIQ_L10C3_FULL_0725
% Mức: VD
\begin{bt}
% ID: AIQ_L10C3_FULL_0725
% Mức: VD
Vật $2\,\text{kg}$ treo đứng yên bằng một sợi dây nhẹ thẳng đứng. Lấy g= $10\,\text{m}$ /s^2. Tính lực căng dây.
\loigiai{
Phân tích lực, chọn hệ quy chiếu và áp dụng định luật hoặc quy tắc phù hợp. Vật cân bằng nên T=P=mg= $20\,\text{N}$.
}
\end{bt}

% ===== Câu 78 =====
% ID: AIQ_L10C3_FULL_0726
% Mức: VDC
\begin{bt}
% ID: AIQ_L10C3_FULL_0726
% Mức: VDC
Vật $3\,\text{kg}$ treo đứng yên bằng một sợi dây nhẹ thẳng đứng. Lấy g= $10\,\text{m}$ /s^2. Tính lực căng dây.
\loigiai{
Phân tích lực, chọn hệ quy chiếu và áp dụng định luật hoặc quy tắc phù hợp. Vật cân bằng nên T=P=mg= $30\,\text{N}$.
}
\end{bt}

% ===== Câu 79 =====
% ID: AIQ_L10C3_FULL_0727
% Mức: VDC
\begin{bt}
% ID: AIQ_L10C3_FULL_0727
% Mức: VDC
Vật $4\,\text{kg}$ treo đứng yên bằng một sợi dây nhẹ thẳng đứng. Lấy g= $10\,\text{m}$ /s^2. Tính lực căng dây.
\loigiai{
Phân tích lực, chọn hệ quy chiếu và áp dụng định luật hoặc quy tắc phù hợp. Vật cân bằng nên T=P=mg= $40\,\text{N}$.
}
\end{bt}

% ===== Câu 80 =====
% ID: AIQ_L10C3_FULL_0728
% Mức: VDC
\begin{bt}
% ID: AIQ_L10C3_FULL_0728
% Mức: VDC
Vật $5\,\text{kg}$ treo đứng yên bằng một sợi dây nhẹ thẳng đứng. Lấy g= $10\,\text{m}$ /s^2. Tính lực căng dây.
\loigiai{
Phân tích lực, chọn hệ quy chiếu và áp dụng định luật hoặc quy tắc phù hợp. Vật cân bằng nên T=P=mg= $50\,\text{N}$.
}
\end{bt}

% ===== Câu 81 =====
% ID: AIQ_L10C3_FULL_0729
% Mức: VDC
\begin{bt}
% ID: AIQ_L10C3_FULL_0729
% Mức: VDC
Vật $6\,\text{kg}$ treo đứng yên bằng một sợi dây nhẹ thẳng đứng. Lấy g= $10\,\text{m}$ /s^2. Tính lực căng dây.
\loigiai{
Phân tích lực, chọn hệ quy chiếu và áp dụng định luật hoặc quy tắc phù hợp. Vật cân bằng nên T=P=mg= $60\,\text{N}$.
}
\end{bt}

% ===== Câu 82 =====
% ID: AIQ_L10C3_FULL_0730
% Mức: NB
\dangbt{Bài toán vật chuyển động có lực căng dây}
\begin{ex}
% ID: AIQ_L10C3_FULL_0730
% Mức: NB
Vật $2\,\text{kg}$ được kéo thẳng đứng lên nhanh dần đều với gia tốc $4\,\text{m}$ /s^2 bằng dây nhẹ. Lấy g= $10\,\text{m}$ /s^2. Tính lực căng dây.
\choice
{28 $N$}
{29 $N$}
{27 $N$}
{\True 28 $N$}
\loigiai{
Chọn D. T−mg=ma nên T=m(g+a)=2(10+4)= $28\,\text{N}$.
}
\end{ex}

% ===== Câu 83 =====
% ID: AIQ_L10C3_FULL_0731
% Mức: NB
\begin{ex}
% ID: AIQ_L10C3_FULL_0731
% Mức: NB
Vật $3\,\text{kg}$ được kéo thẳng đứng lên nhanh dần đều với gia tốc $1\,\text{m}$ /s^2 bằng dây nhẹ. Lấy g= $10\,\text{m}$ /s^2. Tính lực căng dây.
\choice
{\True 33 $N$}
{34 $N$}
{32 $N$}
{66 $N$}
\loigiai{
Chọn A. T−mg=ma nên T=m(g+a)=3(10+1)= $33\,\text{N}$.
}
\end{ex}

% ===== Câu 84 =====
% ID: AIQ_L10C3_FULL_0732
% Mức: NB
\begin{ex}
% ID: AIQ_L10C3_FULL_0732
% Mức: NB
Vật $4\,\text{kg}$ được kéo thẳng đứng lên nhanh dần đều với gia tốc $2\,\text{m}$ /s^2 bằng dây nhẹ. Lấy g= $10\,\text{m}$ /s^2. Tính lực căng dây.
\choice
{48 $N$}
{\True 48 $N$}
{47 $N$}
{96 $N$}
\loigiai{
Chọn B. T−mg=ma nên T=m(g+a)=4(10+2)= $48\,\text{N}$.
}
\end{ex}

% ===== Câu 85 =====
% ID: AIQ_L10C3_FULL_0733
% Mức: NB
\begin{ex}
% ID: AIQ_L10C3_FULL_0733
% Mức: NB
Vật $5\,\text{kg}$ được kéo thẳng đứng lên nhanh dần đều với gia tốc $3\,\text{m}$ /s^2 bằng dây nhẹ. Lấy g= $10\,\text{m}$ /s^2. Tính lực căng dây.
\choice
{65 $N$}
{66 $N$}
{\True 65 $N$}
{130 $N$}
\loigiai{
Chọn C. T−mg=ma nên T=m(g+a)=5(10+3)= $65\,\text{N}$.
}
\end{ex}

% ===== Câu 86 =====
% ID: AIQ_L10C3_FULL_0734
% Mức: TH
\begin{ex}
% ID: AIQ_L10C3_FULL_0734
% Mức: TH
Vật $6\,\text{kg}$ được kéo thẳng đứng lên nhanh dần đều với gia tốc $4\,\text{m}$ /s^2 bằng dây nhẹ. Lấy g= $10\,\text{m}$ /s^2. Tính lực căng dây.
\choice
{84 $N$}
{85 $N$}
{83 $N$}
{\True 84 $N$}
\loigiai{
Chọn D. T−mg=ma nên T=m(g+a)=6(10+4)= $84\,\text{N}$.
}
\end{ex}

% ===== Câu 87 =====
% ID: AIQ_L10C3_FULL_0735
% Mức: TH
\begin{ex}
% ID: AIQ_L10C3_FULL_0735
% Mức: TH
Vật $2\,\text{kg}$ được kéo thẳng đứng lên nhanh dần đều với gia tốc $1\,\text{m}$ /s^2 bằng dây nhẹ. Lấy g= $10\,\text{m}$ /s^2. Tính lực căng dây.
\choice
{\True 22 $N$}
{23 $N$}
{21 $N$}
{44 $N$}
\loigiai{
Chọn A. T−mg=ma nên T=m(g+a)=2(10+1)= $22\,\text{N}$.
}
\end{ex}

% ===== Câu 88 =====
% ID: AIQ_L10C3_FULL_0736
% Mức: TH
\begin{ex}
% ID: AIQ_L10C3_FULL_0736
% Mức: TH
Vật $3\,\text{kg}$ được kéo thẳng đứng lên nhanh dần đều với gia tốc $2\,\text{m}$ /s^2 bằng dây nhẹ. Lấy g= $10\,\text{m}$ /s^2. Tính lực căng dây.
\choice
{36 $N$}
{\True 36 $N$}
{35 $N$}
{72 $N$}
\loigiai{
Chọn B. T−mg=ma nên T=m(g+a)=3(10+2)= $36\,\text{N}$.
}
\end{ex}

% ===== Câu 89 =====
% ID: AIQ_L10C3_FULL_0737
% Mức: TH
\begin{ex}
% ID: AIQ_L10C3_FULL_0737
% Mức: TH
Vật $4\,\text{kg}$ được kéo thẳng đứng lên nhanh dần đều với gia tốc $3\,\text{m}$ /s^2 bằng dây nhẹ. Lấy g= $10\,\text{m}$ /s^2. Tính lực căng dây.
\choice
{52 $N$}
{53 $N$}
{\True 52 $N$}
{104 $N$}
\loigiai{
Chọn C. T−mg=ma nên T=m(g+a)=4(10+3)= $52\,\text{N}$.
}
\end{ex}

% ===== Câu 90 =====
% ID: AIQ_L10C3_FULL_0738
% Mức: VD
\begin{ex}
% ID: AIQ_L10C3_FULL_0738
% Mức: VD
Vật $5\,\text{kg}$ được kéo thẳng đứng lên nhanh dần đều với gia tốc $4\,\text{m}$ /s^2 bằng dây nhẹ. Lấy g= $10\,\text{m}$ /s^2. Tính lực căng dây.
\choice
{70 $N$}
{71 $N$}
{69 $N$}
{\True 70 $N$}
\loigiai{
Chọn D. T−mg=ma nên T=m(g+a)=5(10+4)= $70\,\text{N}$.
}
\end{ex}

% ===== Câu 91 =====
% ID: AIQ_L10C3_FULL_0739
% Mức: VD
\begin{ex}
% ID: AIQ_L10C3_FULL_0739
% Mức: VD
Vật $6\,\text{kg}$ được kéo thẳng đứng lên nhanh dần đều với gia tốc $1\,\text{m}$ /s^2 bằng dây nhẹ. Lấy g= $10\,\text{m}$ /s^2. Tính lực căng dây.
\choice
{\True 66 $N$}
{67 $N$}
{65 $N$}
{132 $N$}
\loigiai{
Chọn A. T−mg=ma nên T=m(g+a)=6(10+1)= $66\,\text{N}$.
}
\end{ex}

% ===== Câu 92 =====
% ID: AIQ_L10C3_FULL_0740
% Mức: TH
\begin{ex}
% ID: AIQ_L10C3_FULL_0740
% Mức: TH
Vật $5\,\text{kg}$ được kéo thẳng đứng lên nhanh dần đều với gia tốc $1\,\text{m}$ /s^2 bằng dây nhẹ. Lấy g= $10\,\text{m}$ /s^2. Tính lực căng dây.
\choiceTF
{\True Giá trị cần tìm là $55\,\text{N}$.}
{Giá trị cần tìm là $56\,\text{N}$.}
{\True Phải xác định đúng các lực tác dụng và chọn chiều dương trước khi lập phương trình.}
{Có thể bỏ qua điều kiện áp dụng và chiều của lực mà kết quả vẫn luôn đúng.}
\loigiai{
\begin{itemchoice}
\item (Đúng) Giá trị cần tìm là $55\,\text{N}$. (Vì): T−mg=ma nên T=m(g+a)=5(10+1)= $55\,\text{N}$. b) Sai vì sai giá trị. c) Đúng. d) Sai vì phải kiểm tra điều kiện áp dụng và chiều lực
\item (Sai) Giá trị cần tìm là $56\,\text{N}$.
\item (Đúng) Phải xác định đúng các lực tác dụng và chọn chiều dương trước khi lập phương trình.
\item (Sai) Có thể bỏ qua điều kiện áp dụng và chiều của lực mà kết quả vẫn luôn đúng.
\end{itemchoice}
}
\end{ex}

% ===== Câu 93 =====
% ID: AIQ_L10C3_FULL_0741
% Mức: TH
\begin{ex}
% ID: AIQ_L10C3_FULL_0741
% Mức: TH
Vật $6\,\text{kg}$ được kéo thẳng đứng lên nhanh dần đều với gia tốc $2\,\text{m}$ /s^2 bằng dây nhẹ. Lấy g= $10\,\text{m}$ /s^2. Tính lực căng dây.
\choiceTF
{\True Giá trị cần tìm là $72\,\text{N}$.}
{Giá trị cần tìm là $73\,\text{N}$.}
{\True Phải xác định đúng các lực tác dụng và chọn chiều dương trước khi lập phương trình.}
{Có thể bỏ qua điều kiện áp dụng và chiều của lực mà kết quả vẫn luôn đúng.}
\loigiai{
\begin{itemchoice}
\item (Đúng) Giá trị cần tìm là $72\,\text{N}$. (Vì): T−mg=ma nên T=m(g+a)=6(10+2)= $72\,\text{N}$. b) Sai vì sai giá trị. c) Đúng. d) Sai vì phải kiểm tra điều kiện áp dụng và chiều lực
\item (Sai) Giá trị cần tìm là $73\,\text{N}$.
\item (Đúng) Phải xác định đúng các lực tác dụng và chọn chiều dương trước khi lập phương trình.
\item (Sai) Có thể bỏ qua điều kiện áp dụng và chiều của lực mà kết quả vẫn luôn đúng.
\end{itemchoice}
}
\end{ex}

% ===== Câu 94 =====
% ID: AIQ_L10C3_FULL_0742
% Mức: VD
\begin{ex}
% ID: AIQ_L10C3_FULL_0742
% Mức: VD
Vật $2\,\text{kg}$ được kéo thẳng đứng lên nhanh dần đều với gia tốc $3\,\text{m}$ /s^2 bằng dây nhẹ. Lấy g= $10\,\text{m}$ /s^2. Tính lực căng dây.
\choiceTF
{\True Giá trị cần tìm là $26\,\text{N}$.}
{Giá trị cần tìm là $27\,\text{N}$.}
{\True Phải xác định đúng các lực tác dụng và chọn chiều dương trước khi lập phương trình.}
{Có thể bỏ qua điều kiện áp dụng và chiều của lực mà kết quả vẫn luôn đúng.}
\loigiai{
\begin{itemchoice}
\item (Đúng) Giá trị cần tìm là $26\,\text{N}$. (Vì): T−mg=ma nên T=m(g+a)=2(10+3)= $26\,\text{N}$. b) Sai vì sai giá trị. c) Đúng. d) Sai vì phải kiểm tra điều kiện áp dụng và chiều lực
\item (Sai) Giá trị cần tìm là $27\,\text{N}$.
\item (Đúng) Phải xác định đúng các lực tác dụng và chọn chiều dương trước khi lập phương trình.
\item (Sai) Có thể bỏ qua điều kiện áp dụng và chiều của lực mà kết quả vẫn luôn đúng.
\end{itemchoice}
}
\end{ex}

% ===== Câu 95 =====
% ID: AIQ_L10C3_FULL_0743
% Mức: VD
\begin{ex}
% ID: AIQ_L10C3_FULL_0743
% Mức: VD
Vật $3\,\text{kg}$ được kéo thẳng đứng lên nhanh dần đều với gia tốc $4\,\text{m}$ /s^2 bằng dây nhẹ. Lấy g= $10\,\text{m}$ /s^2. Tính lực căng dây.
\choiceTF
{\True Giá trị cần tìm là $42\,\text{N}$.}
{Giá trị cần tìm là $43\,\text{N}$.}
{\True Phải xác định đúng các lực tác dụng và chọn chiều dương trước khi lập phương trình.}
{Có thể bỏ qua điều kiện áp dụng và chiều của lực mà kết quả vẫn luôn đúng.}
\loigiai{
\begin{itemchoice}
\item (Đúng) Giá trị cần tìm là $42\,\text{N}$. (Vì): T−mg=ma nên T=m(g+a)=3(10+4)= $42\,\text{N}$. b) Sai vì sai giá trị. c) Đúng. d) Sai vì phải kiểm tra điều kiện áp dụng và chiều lực
\item (Sai) Giá trị cần tìm là $43\,\text{N}$.
\item (Đúng) Phải xác định đúng các lực tác dụng và chọn chiều dương trước khi lập phương trình.
\item (Sai) Có thể bỏ qua điều kiện áp dụng và chiều của lực mà kết quả vẫn luôn đúng.
\end{itemchoice}
}
\end{ex}

% ===== Câu 96 =====
% ID: AIQ_L10C3_FULL_0744
% Mức: VD
\begin{ex}
% ID: AIQ_L10C3_FULL_0744
% Mức: VD
Vật $4\,\text{kg}$ được kéo thẳng đứng lên nhanh dần đều với gia tốc $1\,\text{m}$ /s^2 bằng dây nhẹ. Lấy g= $10\,\text{m}$ /s^2. Tính lực căng dây.
\choiceTF
{\True Giá trị cần tìm là $44\,\text{N}$.}
{Giá trị cần tìm là $45\,\text{N}$.}
{\True Phải xác định đúng các lực tác dụng và chọn chiều dương trước khi lập phương trình.}
{Có thể bỏ qua điều kiện áp dụng và chiều của lực mà kết quả vẫn luôn đúng.}
\loigiai{
\begin{itemchoice}
\item (Đúng) Giá trị cần tìm là $44\,\text{N}$. (Vì): T−mg=ma nên T=m(g+a)=4(10+1)= $44\,\text{N}$. b) Sai vì sai giá trị. c) Đúng. d) Sai vì phải kiểm tra điều kiện áp dụng và chiều lực
\item (Sai) Giá trị cần tìm là $45\,\text{N}$.
\item (Đúng) Phải xác định đúng các lực tác dụng và chọn chiều dương trước khi lập phương trình.
\item (Sai) Có thể bỏ qua điều kiện áp dụng và chiều của lực mà kết quả vẫn luôn đúng.
\end{itemchoice}
}
\end{ex}

% ===== Câu 97 =====
% ID: AIQ_L10C3_FULL_0745
% Mức: TH
\begin{ex}
% ID: AIQ_L10C3_FULL_0745
% Mức: TH
Vật $3\,\text{kg}$ được kéo thẳng đứng lên nhanh dần đều với gia tốc $2\,\text{m}$ /s^2 bằng dây nhẹ. Lấy g= $10\,\text{m}$ /s^2. Tính lực căng dây.
\shortans{36}
\loigiai{
T−mg=ma nên T=m(g+a)=3(10+2)= $36\,\text{N}$.
}
\end{ex}

% ===== Câu 98 =====
% ID: AIQ_L10C3_FULL_0746
% Mức: TH
\begin{ex}
% ID: AIQ_L10C3_FULL_0746
% Mức: TH
Vật $4\,\text{kg}$ được kéo thẳng đứng lên nhanh dần đều với gia tốc $3\,\text{m}$ /s^2 bằng dây nhẹ. Lấy g= $10\,\text{m}$ /s^2. Tính lực căng dây.
\shortans{52}
\loigiai{
T−mg=ma nên T=m(g+a)=4(10+3)= $52\,\text{N}$.
}
\end{ex}

% ===== Câu 99 =====
% ID: AIQ_L10C3_FULL_0747
% Mức: VD
\begin{ex}
% ID: AIQ_L10C3_FULL_0747
% Mức: VD
Vật $5\,\text{kg}$ được kéo thẳng đứng lên nhanh dần đều với gia tốc $4\,\text{m}$ /s^2 bằng dây nhẹ. Lấy g= $10\,\text{m}$ /s^2. Tính lực căng dây.
\shortans{70}
\loigiai{
T−mg=ma nên T=m(g+a)=5(10+4)= $70\,\text{N}$.
}
\end{ex}

% ===== Câu 100 =====
% ID: AIQ_L10C3_FULL_0748
% Mức: VD
\begin{ex}
% ID: AIQ_L10C3_FULL_0748
% Mức: VD
Vật $6\,\text{kg}$ được kéo thẳng đứng lên nhanh dần đều với gia tốc $1\,\text{m}$ /s^2 bằng dây nhẹ. Lấy g= $10\,\text{m}$ /s^2. Tính lực căng dây.
\shortans{66}
\loigiai{
T−mg=ma nên T=m(g+a)=6(10+1)= $66\,\text{N}$.
}
\end{ex}

% ===== Câu 101 =====
% ID: AIQ_L10C3_FULL_0749
% Mức: VD
\begin{ex}
% ID: AIQ_L10C3_FULL_0749
% Mức: VD
Vật $2\,\text{kg}$ được kéo thẳng đứng lên nhanh dần đều với gia tốc $2\,\text{m}$ /s^2 bằng dây nhẹ. Lấy g= $10\,\text{m}$ /s^2. Tính lực căng dây.
\shortans{24}
\loigiai{
T−mg=ma nên T=m(g+a)=2(10+2)= $24\,\text{N}$.
}
\end{ex}

% ===== Câu 102 =====
% ID: AIQ_L10C3_FULL_0750
% Mức: VDC
\begin{ex}
% ID: AIQ_L10C3_FULL_0750
% Mức: VDC
Vật $3\,\text{kg}$ được kéo thẳng đứng lên nhanh dần đều với gia tốc $3\,\text{m}$ /s^2 bằng dây nhẹ. Lấy g= $10\,\text{m}$ /s^2. Tính lực căng dây.
\shortans{39}
\loigiai{
T−mg=ma nên T=m(g+a)=3(10+3)= $39\,\text{N}$.
}
\end{ex}

% ===== Câu 103 =====
% ID: AIQ_L10C3_FULL_0751
% Mức: VD
\begin{bt}
% ID: AIQ_L10C3_FULL_0751
% Mức: VD
Vật $6\,\text{kg}$ được kéo thẳng đứng lên nhanh dần đều với gia tốc $3\,\text{m}$ /s^2 bằng dây nhẹ. Lấy g= $10\,\text{m}$ /s^2. Tính lực căng dây.
\loigiai{
Phân tích lực, chọn hệ quy chiếu và áp dụng định luật hoặc quy tắc phù hợp. T−mg=ma nên T=m(g+a)=6(10+3)= $78\,\text{N}$.
}
\end{bt}

% ===== Câu 104 =====
% ID: AIQ_L10C3_FULL_0752
% Mức: VD
\begin{bt}
% ID: AIQ_L10C3_FULL_0752
% Mức: VD
Vật $2\,\text{kg}$ được kéo thẳng đứng lên nhanh dần đều với gia tốc $4\,\text{m}$ /s^2 bằng dây nhẹ. Lấy g= $10\,\text{m}$ /s^2. Tính lực căng dây.
\loigiai{
Phân tích lực, chọn hệ quy chiếu và áp dụng định luật hoặc quy tắc phù hợp. T−mg=ma nên T=m(g+a)=2(10+4)= $28\,\text{N}$.
}
\end{bt}

% ===== Câu 105 =====
% ID: AIQ_L10C3_FULL_0753
% Mức: VDC
\begin{bt}
% ID: AIQ_L10C3_FULL_0753
% Mức: VDC
Vật $3\,\text{kg}$ được kéo thẳng đứng lên nhanh dần đều với gia tốc $1\,\text{m}$ /s^2 bằng dây nhẹ. Lấy g= $10\,\text{m}$ /s^2. Tính lực căng dây.
\loigiai{
Phân tích lực, chọn hệ quy chiếu và áp dụng định luật hoặc quy tắc phù hợp. T−mg=ma nên T=m(g+a)=3(10+1)= $33\,\text{N}$.
}
\end{bt}

% ===== Câu 106 =====
% ID: AIQ_L10C3_FULL_0754
% Mức: VDC
\begin{bt}
% ID: AIQ_L10C3_FULL_0754
% Mức: VDC
Vật $4\,\text{kg}$ được kéo thẳng đứng lên nhanh dần đều với gia tốc $2\,\text{m}$ /s^2 bằng dây nhẹ. Lấy g= $10\,\text{m}$ /s^2. Tính lực căng dây.
\loigiai{
Phân tích lực, chọn hệ quy chiếu và áp dụng định luật hoặc quy tắc phù hợp. T−mg=ma nên T=m(g+a)=4(10+2)= $48\,\text{N}$.
}
\end{bt}

% ===== Câu 107 =====
% ID: AIQ_L10C3_FULL_0755
% Mức: VDC
\begin{bt}
% ID: AIQ_L10C3_FULL_0755
% Mức: VDC
Vật $5\,\text{kg}$ được kéo thẳng đứng lên nhanh dần đều với gia tốc $3\,\text{m}$ /s^2 bằng dây nhẹ. Lấy g= $10\,\text{m}$ /s^2. Tính lực căng dây.
\loigiai{
Phân tích lực, chọn hệ quy chiếu và áp dụng định luật hoặc quy tắc phù hợp. T−mg=ma nên T=m(g+a)=5(10+3)= $65\,\text{N}$.
}
\end{bt}

% ===== Câu 108 =====
% ID: AIQ_L10C3_FULL_0756
% Mức: VDC
\begin{bt}
% ID: AIQ_L10C3_FULL_0756
% Mức: VDC
Vật $6\,\text{kg}$ được kéo thẳng đứng lên nhanh dần đều với gia tốc $4\,\text{m}$ /s^2 bằng dây nhẹ. Lấy g= $10\,\text{m}$ /s^2. Tính lực căng dây.
\loigiai{
Phân tích lực, chọn hệ quy chiếu và áp dụng định luật hoặc quy tắc phù hợp. T−mg=ma nên T=m(g+a)=6(10+4)= $84\,\text{N}$.
}
\end{bt}

% ===== Câu 109 =====
% ID: AIQ_L10C3_FULL_0757
% Mức: NB
\dangbt{Hệ vật nối với nhau bằng dây nhẹ}
\begin{ex}
% ID: AIQ_L10C3_FULL_0757
% Mức: NB
Máy Atwood gồm hai vật $2\,\text{kg}$ và $4\,\text{kg}$ nối bằng dây nhẹ qua ròng rọc lí tưởng. Lấy g= $10\,\text{m}$ /s^2. Tính độ lớn gia tốc của hệ.
\choice
{\True $3,33 m/s^2$}
{$4,33 m/s^2$}
{$2,33 m/s^2$}
{$6,66 m/s^2$}
\loigiai{
Chọn A. a=(m2−m1)g/(m1+m2)=20/6= $3,33\,\text{m}$ /s^2.
}
\end{ex}

% ===== Câu 110 =====
% ID: AIQ_L10C3_FULL_0758
% Mức: NB
\begin{ex}
% ID: AIQ_L10C3_FULL_0758
% Mức: NB
Máy Atwood gồm hai vật $3\,\text{kg}$ và $5\,\text{kg}$ nối bằng dây nhẹ qua ròng rọc lí tưởng. Lấy g= $10\,\text{m}$ /s^2. Tính độ lớn gia tốc của hệ.
\choice
{$2,5 m/s^2$}
{\True $2,5 m/s^2$}
{$1,5 m/s^2$}
{$5 m/s^2$}
\loigiai{
Chọn B. a=(m2−m1)g/(m1+m2)=20/8= $2,5\,\text{m}$ /s^2.
}
\end{ex}

% ===== Câu 111 =====
% ID: AIQ_L10C3_FULL_0759
% Mức: NB
\begin{ex}
% ID: AIQ_L10C3_FULL_0759
% Mức: NB
Máy Atwood gồm hai vật $4\,\text{kg}$ và $6\,\text{kg}$ nối bằng dây nhẹ qua ròng rọc lí tưởng. Lấy g= $10\,\text{m}$ /s^2. Tính độ lớn gia tốc của hệ.
\choice
{$2 m/s^2$}
{$3 m/s^2$}
{\True $2 m/s^2$}
{$4 m/s^2$}
\loigiai{
Chọn C. a=(m2−m1)g/(m1+m2)=20/10= $2\,\text{m}$ /s^2.
}
\end{ex}

% ===== Câu 112 =====
% ID: AIQ_L10C3_FULL_0760
% Mức: NB
\begin{ex}
% ID: AIQ_L10C3_FULL_0760
% Mức: NB
Máy Atwood gồm hai vật $5\,\text{kg}$ và $7\,\text{kg}$ nối bằng dây nhẹ qua ròng rọc lí tưởng. Lấy g= $10\,\text{m}$ /s^2. Tính độ lớn gia tốc của hệ.
\choice
{$1,67 m/s^2$}
{$2,67 m/s^2$}
{$0,67 m/s^2$}
{\True $1,67 m/s^2$}
\loigiai{
Chọn D. a=(m2−m1)g/(m1+m2)=20/12= $1,67\,\text{m}$ /s^2.
}
\end{ex}

% ===== Câu 113 =====
% ID: AIQ_L10C3_FULL_0761
% Mức: TH
\begin{ex}
% ID: AIQ_L10C3_FULL_0761
% Mức: TH
Máy Atwood gồm hai vật $6\,\text{kg}$ và $8\,\text{kg}$ nối bằng dây nhẹ qua ròng rọc lí tưởng. Lấy g= $10\,\text{m}$ /s^2. Tính độ lớn gia tốc của hệ.
\choice
{\True $1,43 m/s^2$}
{$2,43 m/s^2$}
{$0,43 m/s^2$}
{$2,86 m/s^2$}
\loigiai{
Chọn A. a=(m2−m1)g/(m1+m2)=20/14= $1,43\,\text{m}$ /s^2.
}
\end{ex}

% ===== Câu 114 =====
% ID: AIQ_L10C3_FULL_0762
% Mức: TH
\begin{ex}
% ID: AIQ_L10C3_FULL_0762
% Mức: TH
Máy Atwood gồm hai vật $2\,\text{kg}$ và $4\,\text{kg}$ nối bằng dây nhẹ qua ròng rọc lí tưởng. Lấy g= $10\,\text{m}$ /s^2. Tính độ lớn gia tốc của hệ.
\choice
{$3,33 m/s^2$}
{\True $3,33 m/s^2$}
{$2,33 m/s^2$}
{$6,66 m/s^2$}
\loigiai{
Chọn B. a=(m2−m1)g/(m1+m2)=20/6= $3,33\,\text{m}$ /s^2.
}
\end{ex}

% ===== Câu 115 =====
% ID: AIQ_L10C3_FULL_0763
% Mức: TH
\begin{ex}
% ID: AIQ_L10C3_FULL_0763
% Mức: TH
Máy Atwood gồm hai vật $3\,\text{kg}$ và $5\,\text{kg}$ nối bằng dây nhẹ qua ròng rọc lí tưởng. Lấy g= $10\,\text{m}$ /s^2. Tính độ lớn gia tốc của hệ.
\choice
{$2,5 m/s^2$}
{$3,5 m/s^2$}
{\True $2,5 m/s^2$}
{$5 m/s^2$}
\loigiai{
Chọn C. a=(m2−m1)g/(m1+m2)=20/8= $2,5\,\text{m}$ /s^2.
}
\end{ex}

% ===== Câu 116 =====
% ID: AIQ_L10C3_FULL_0764
% Mức: TH
\begin{ex}
% ID: AIQ_L10C3_FULL_0764
% Mức: TH
Máy Atwood gồm hai vật $4\,\text{kg}$ và $6\,\text{kg}$ nối bằng dây nhẹ qua ròng rọc lí tưởng. Lấy g= $10\,\text{m}$ /s^2. Tính độ lớn gia tốc của hệ.
\choice
{$2 m/s^2$}
{$3 m/s^2$}
{$1 m/s^2$}
{\True $2 m/s^2$}
\loigiai{
Chọn D. a=(m2−m1)g/(m1+m2)=20/10= $2\,\text{m}$ /s^2.
}
\end{ex}

% ===== Câu 117 =====
% ID: AIQ_L10C3_FULL_0765
% Mức: VD
\begin{ex}
% ID: AIQ_L10C3_FULL_0765
% Mức: VD
Máy Atwood gồm hai vật $5\,\text{kg}$ và $7\,\text{kg}$ nối bằng dây nhẹ qua ròng rọc lí tưởng. Lấy g= $10\,\text{m}$ /s^2. Tính độ lớn gia tốc của hệ.
\choice
{\True $1,67 m/s^2$}
{$2,67 m/s^2$}
{$0,67 m/s^2$}
{$3,34 m/s^2$}
\loigiai{
Chọn A. a=(m2−m1)g/(m1+m2)=20/12= $1,67\,\text{m}$ /s^2.
}
\end{ex}

% ===== Câu 118 =====
% ID: AIQ_L10C3_FULL_0766
% Mức: VD
\begin{ex}
% ID: AIQ_L10C3_FULL_0766
% Mức: VD
Máy Atwood gồm hai vật $6\,\text{kg}$ và $8\,\text{kg}$ nối bằng dây nhẹ qua ròng rọc lí tưởng. Lấy g= $10\,\text{m}$ /s^2. Tính độ lớn gia tốc của hệ.
\choice
{$1,43 m/s^2$}
{\True $1,43 m/s^2$}
{$0,43 m/s^2$}
{$2,86 m/s^2$}
\loigiai{
Chọn B. a=(m2−m1)g/(m1+m2)=20/14= $1,43\,\text{m}$ /s^2.
}
\end{ex}

% ===== Câu 119 =====
% ID: AIQ_L10C3_FULL_0767
% Mức: TH
\begin{ex}
% ID: AIQ_L10C3_FULL_0767
% Mức: TH
Máy Atwood gồm hai vật $5\,\text{kg}$ và $7\,\text{kg}$ nối bằng dây nhẹ qua ròng rọc lí tưởng. Lấy g= $10\,\text{m}$ /s^2. Tính độ lớn gia tốc của hệ.
\choiceTF
{\True Giá trị cần tìm là $1,67\,\text{m}$ /s^2.}
{Giá trị cần tìm là $2,67\,\text{m}$ /s^2.}
{\True Phải xác định đúng các lực tác dụng và chọn chiều dương trước khi lập phương trình.}
{Có thể bỏ qua điều kiện áp dụng và chiều của lực mà kết quả vẫn luôn đúng.}
\loigiai{
\begin{itemchoice}
\item (Đúng) Giá trị cần tìm là $1,67\,\text{m}$ /s^2. (Vì): a=(m2−m1)g/(m1+m2)=20/12= $1,67\,\text{m}$ /s^2. b) Sai vì sai giá trị. c) Đúng. d) Sai vì phải kiểm tra điều kiện áp dụng và chiều lực
\item (Sai) Giá trị cần tìm là $2,67\,\text{m}$ /s^2.
\item (Đúng) Phải xác định đúng các lực tác dụng và chọn chiều dương trước khi lập phương trình.
\item (Sai) Có thể bỏ qua điều kiện áp dụng và chiều của lực mà kết quả vẫn luôn đúng.
\end{itemchoice}
}
\end{ex}

% ===== Câu 120 =====
% ID: AIQ_L10C3_FULL_0768
% Mức: TH
\begin{ex}
% ID: AIQ_L10C3_FULL_0768
% Mức: TH
Máy Atwood gồm hai vật $6\,\text{kg}$ và $8\,\text{kg}$ nối bằng dây nhẹ qua ròng rọc lí tưởng. Lấy g= $10\,\text{m}$ /s^2. Tính độ lớn gia tốc của hệ.
\choiceTF
{\True Giá trị cần tìm là $1,43\,\text{m}$ /s^2.}
{Giá trị cần tìm là $2,43\,\text{m}$ /s^2.}
{\True Phải xác định đúng các lực tác dụng và chọn chiều dương trước khi lập phương trình.}
{Có thể bỏ qua điều kiện áp dụng và chiều của lực mà kết quả vẫn luôn đúng.}
\loigiai{
\begin{itemchoice}
\item (Đúng) Giá trị cần tìm là $1,43\,\text{m}$ /s^2. (Vì): a=(m2−m1)g/(m1+m2)=20/14= $1,43\,\text{m}$ /s^2. b) Sai vì sai giá trị. c) Đúng. d) Sai vì phải kiểm tra điều kiện áp dụng và chiều lực
\item (Sai) Giá trị cần tìm là $2,43\,\text{m}$ /s^2.
\item (Đúng) Phải xác định đúng các lực tác dụng và chọn chiều dương trước khi lập phương trình.
\item (Sai) Có thể bỏ qua điều kiện áp dụng và chiều của lực mà kết quả vẫn luôn đúng.
\end{itemchoice}
}
\end{ex}

% ===== Câu 121 =====
% ID: AIQ_L10C3_FULL_0769
% Mức: VD
\begin{ex}
% ID: AIQ_L10C3_FULL_0769
% Mức: VD
Máy Atwood gồm hai vật $2\,\text{kg}$ và $4\,\text{kg}$ nối bằng dây nhẹ qua ròng rọc lí tưởng. Lấy g= $10\,\text{m}$ /s^2. Tính độ lớn gia tốc của hệ.
\choiceTF
{\True Giá trị cần tìm là $3,33\,\text{m}$ /s^2.}
{Giá trị cần tìm là $4,33\,\text{m}$ /s^2.}
{\True Phải xác định đúng các lực tác dụng và chọn chiều dương trước khi lập phương trình.}
{Có thể bỏ qua điều kiện áp dụng và chiều của lực mà kết quả vẫn luôn đúng.}
\loigiai{
\begin{itemchoice}
\item (Đúng) Giá trị cần tìm là $3,33\,\text{m}$ /s^2. (Vì): a=(m2−m1)g/(m1+m2)=20/6= $3,33\,\text{m}$ /s^2. b) Sai vì sai giá trị. c) Đúng. d) Sai vì phải kiểm tra điều kiện áp dụng và chiều lực
\item (Sai) Giá trị cần tìm là $4,33\,\text{m}$ /s^2.
\item (Đúng) Phải xác định đúng các lực tác dụng và chọn chiều dương trước khi lập phương trình.
\item (Sai) Có thể bỏ qua điều kiện áp dụng và chiều của lực mà kết quả vẫn luôn đúng.
\end{itemchoice}
}
\end{ex}

% ===== Câu 122 =====
% ID: AIQ_L10C3_FULL_0770
% Mức: VD
\begin{ex}
% ID: AIQ_L10C3_FULL_0770
% Mức: VD
Máy Atwood gồm hai vật $3\,\text{kg}$ và $5\,\text{kg}$ nối bằng dây nhẹ qua ròng rọc lí tưởng. Lấy g= $10\,\text{m}$ /s^2. Tính độ lớn gia tốc của hệ.
\choiceTF
{\True Giá trị cần tìm là $2,5\,\text{m}$ /s^2.}
{Giá trị cần tìm là $3,5\,\text{m}$ /s^2.}
{\True Phải xác định đúng các lực tác dụng và chọn chiều dương trước khi lập phương trình.}
{Có thể bỏ qua điều kiện áp dụng và chiều của lực mà kết quả vẫn luôn đúng.}
\loigiai{
\begin{itemchoice}
\item (Đúng) Giá trị cần tìm là $2,5\,\text{m}$ /s^2. (Vì): a=(m2−m1)g/(m1+m2)=20/8= $2,5\,\text{m}$ /s^2. b) Sai vì sai giá trị. c) Đúng. d) Sai vì phải kiểm tra điều kiện áp dụng và chiều lực
\item (Sai) Giá trị cần tìm là $3,5\,\text{m}$ /s^2.
\item (Đúng) Phải xác định đúng các lực tác dụng và chọn chiều dương trước khi lập phương trình.
\item (Sai) Có thể bỏ qua điều kiện áp dụng và chiều của lực mà kết quả vẫn luôn đúng.
\end{itemchoice}
}
\end{ex}

% ===== Câu 123 =====
% ID: AIQ_L10C3_FULL_0771
% Mức: VD
\begin{ex}
% ID: AIQ_L10C3_FULL_0771
% Mức: VD
Máy Atwood gồm hai vật $4\,\text{kg}$ và $6\,\text{kg}$ nối bằng dây nhẹ qua ròng rọc lí tưởng. Lấy g= $10\,\text{m}$ /s^2. Tính độ lớn gia tốc của hệ.
\choiceTF
{\True Giá trị cần tìm là $2\,\text{m}$ /s^2.}
{Giá trị cần tìm là $3\,\text{m}$ /s^2.}
{\True Phải xác định đúng các lực tác dụng và chọn chiều dương trước khi lập phương trình.}
{Có thể bỏ qua điều kiện áp dụng và chiều của lực mà kết quả vẫn luôn đúng.}
\loigiai{
\begin{itemchoice}
\item (Đúng) Giá trị cần tìm là $2\,\text{m}$ /s^2. (Vì): a=(m2−m1)g/(m1+m2)=20/10= $2\,\text{m}$ /s^2. b) Sai vì sai giá trị. c) Đúng. d) Sai vì phải kiểm tra điều kiện áp dụng và chiều lực
\item (Sai) Giá trị cần tìm là $3\,\text{m}$ /s^2.
\item (Đúng) Phải xác định đúng các lực tác dụng và chọn chiều dương trước khi lập phương trình.
\item (Sai) Có thể bỏ qua điều kiện áp dụng và chiều của lực mà kết quả vẫn luôn đúng.
\end{itemchoice}
}
\end{ex}

% ===== Câu 124 =====
% ID: AIQ_L10C3_FULL_0772
% Mức: TH
\begin{ex}
% ID: AIQ_L10C3_FULL_0772
% Mức: TH
Máy Atwood gồm hai vật $3\,\text{kg}$ và $5\,\text{kg}$ nối bằng dây nhẹ qua ròng rọc lí tưởng. Lấy g= $10\,\text{m}$ /s^2. Tính độ lớn gia tốc của hệ.
\shortans{2,5}
\loigiai{
$a=(m2−m1)g/(m1+m2)=20/8=2,5 m/s^2.$
}
\end{ex}

% ===== Câu 125 =====
% ID: AIQ_L10C3_FULL_0773
% Mức: TH
\begin{ex}
% ID: AIQ_L10C3_FULL_0773
% Mức: TH
Máy Atwood gồm hai vật $4\,\text{kg}$ và $6\,\text{kg}$ nối bằng dây nhẹ qua ròng rọc lí tưởng. Lấy g= $10\,\text{m}$ /s^2. Tính độ lớn gia tốc của hệ.
\shortans{2}
\loigiai{
$a=(m2−m1)g/(m1+m2)=20/10=2 m/s^2.$
}
\end{ex}

% ===== Câu 126 =====
% ID: AIQ_L10C3_FULL_0774
% Mức: VD
\begin{ex}
% ID: AIQ_L10C3_FULL_0774
% Mức: VD
Máy Atwood gồm hai vật $5\,\text{kg}$ và $7\,\text{kg}$ nối bằng dây nhẹ qua ròng rọc lí tưởng. Lấy g= $10\,\text{m}$ /s^2. Tính độ lớn gia tốc của hệ.
\shortans{1,67}
\loigiai{
$a=(m2−m1)g/(m1+m2)=20/12=1,67 m/s^2.$
}
\end{ex}

% ===== Câu 127 =====
% ID: AIQ_L10C3_FULL_0775
% Mức: VD
\begin{ex}
% ID: AIQ_L10C3_FULL_0775
% Mức: VD
Máy Atwood gồm hai vật $6\,\text{kg}$ và $8\,\text{kg}$ nối bằng dây nhẹ qua ròng rọc lí tưởng. Lấy g= $10\,\text{m}$ /s^2. Tính độ lớn gia tốc của hệ.
\shortans{1,43}
\loigiai{
$a=(m2−m1)g/(m1+m2)=20/14=1,43 m/s^2.$
}
\end{ex}

% ===== Câu 128 =====
% ID: AIQ_L10C3_FULL_0776
% Mức: VD
\begin{ex}
% ID: AIQ_L10C3_FULL_0776
% Mức: VD
Máy Atwood gồm hai vật $2\,\text{kg}$ và $4\,\text{kg}$ nối bằng dây nhẹ qua ròng rọc lí tưởng. Lấy g= $10\,\text{m}$ /s^2. Tính độ lớn gia tốc của hệ.
\shortans{3,33}
\loigiai{
$a=(m2−m1)g/(m1+m2)=20/6=3,33 m/s^2.$
}
\end{ex}

% ===== Câu 129 =====
% ID: AIQ_L10C3_FULL_0777
% Mức: VDC
\begin{ex}
% ID: AIQ_L10C3_FULL_0777
% Mức: VDC
Máy Atwood gồm hai vật $3\,\text{kg}$ và $5\,\text{kg}$ nối bằng dây nhẹ qua ròng rọc lí tưởng. Lấy g= $10\,\text{m}$ /s^2. Tính độ lớn gia tốc của hệ.
\shortans{2,5}
\loigiai{
$a=(m2−m1)g/(m1+m2)=20/8=2,5 m/s^2.$
}
\end{ex}

% ===== Câu 130 =====
% ID: AIQ_L10C3_FULL_0778
% Mức: VD
\begin{bt}
% ID: AIQ_L10C3_FULL_0778
% Mức: VD
Máy Atwood gồm hai vật $6\,\text{kg}$ và $8\,\text{kg}$ nối bằng dây nhẹ qua ròng rọc lí tưởng. Lấy g= $10\,\text{m}$ /s^2. Tính độ lớn gia tốc của hệ.
\loigiai{
Phân tích lực, chọn hệ quy chiếu và áp dụng định luật hoặc quy tắc phù hợp. a=(m2−m1)g/(m1+m2)=20/14= $1,43\,\text{m}$ /s^2.
}
\end{bt}

% ===== Câu 131 =====
% ID: AIQ_L10C3_FULL_0779
% Mức: VD
\begin{bt}
% ID: AIQ_L10C3_FULL_0779
% Mức: VD
Máy Atwood gồm hai vật $2\,\text{kg}$ và $4\,\text{kg}$ nối bằng dây nhẹ qua ròng rọc lí tưởng. Lấy g= $10\,\text{m}$ /s^2. Tính độ lớn gia tốc của hệ.
\loigiai{
Phân tích lực, chọn hệ quy chiếu và áp dụng định luật hoặc quy tắc phù hợp. a=(m2−m1)g/(m1+m2)=20/6= $3,33\,\text{m}$ /s^2.
}
\end{bt}

% ===== Câu 132 =====
% ID: AIQ_L10C3_FULL_0780
% Mức: VDC
\begin{bt}
% ID: AIQ_L10C3_FULL_0780
% Mức: VDC
Máy Atwood gồm hai vật $3\,\text{kg}$ và $5\,\text{kg}$ nối bằng dây nhẹ qua ròng rọc lí tưởng. Lấy g= $10\,\text{m}$ /s^2. Tính độ lớn gia tốc của hệ.
\loigiai{
Phân tích lực, chọn hệ quy chiếu và áp dụng định luật hoặc quy tắc phù hợp. a=(m2−m1)g/(m1+m2)=20/8= $2,5\,\text{m}$ /s^2.
}
\end{bt}

% ===== Câu 133 =====
% ID: AIQ_L10C3_FULL_0781
% Mức: VDC
\begin{bt}
% ID: AIQ_L10C3_FULL_0781
% Mức: VDC
Máy Atwood gồm hai vật $4\,\text{kg}$ và $6\,\text{kg}$ nối bằng dây nhẹ qua ròng rọc lí tưởng. Lấy g= $10\,\text{m}$ /s^2. Tính độ lớn gia tốc của hệ.
\loigiai{
Phân tích lực, chọn hệ quy chiếu và áp dụng định luật hoặc quy tắc phù hợp. a=(m2−m1)g/(m1+m2)=20/10= $2\,\text{m}$ /s^2.
}
\end{bt}

% ===== Câu 134 =====
% ID: AIQ_L10C3_FULL_0782
% Mức: VDC
\begin{bt}
% ID: AIQ_L10C3_FULL_0782
% Mức: VDC
Máy Atwood gồm hai vật $5\,\text{kg}$ và $7\,\text{kg}$ nối bằng dây nhẹ qua ròng rọc lí tưởng. Lấy g= $10\,\text{m}$ /s^2. Tính độ lớn gia tốc của hệ.
\loigiai{
Phân tích lực, chọn hệ quy chiếu và áp dụng định luật hoặc quy tắc phù hợp. a=(m2−m1)g/(m1+m2)=20/12= $1,67\,\text{m}$ /s^2.
}
\end{bt}

% ===== Câu 135 =====
% ID: AIQ_L10C3_FULL_0783
% Mức: VDC
\begin{bt}
% ID: AIQ_L10C3_FULL_0783
% Mức: VDC
Máy Atwood gồm hai vật $6\,\text{kg}$ và $8\,\text{kg}$ nối bằng dây nhẹ qua ròng rọc lí tưởng. Lấy g= $10\,\text{m}$ /s^2. Tính độ lớn gia tốc của hệ.
\loigiai{
Phân tích lực, chọn hệ quy chiếu và áp dụng định luật hoặc quy tắc phù hợp. a=(m2−m1)g/(m1+m2)=20/14= $1,43\,\text{m}$ /s^2.
}
\end{bt}

% ===== Câu 136 =====
% ID: AIQ_L10C3_FULL_0784
% Mức: NB
\dangbt{Bài toán tổng hợp trọng lực và lực căng trong thực tiễn}
\begin{ex}
% ID: AIQ_L10C3_FULL_0784
% Mức: NB
Biển quảng cáo khối lượng $2\,\text{kg}$ được treo cân bằng bằng hai dây giống nhau, đối xứng; mỗi dây hợp với phương ngang góc 30°. Lấy g= $10\,\text{m}$ /s^2. Tính lực căng mỗi dây.
\choice
{20 $N$}
{\True 20 $N$}
{19 $N$}
{40 $N$}
\loigiai{
Chọn B. 2 $T$ sin30°=mg nên T=mg= $20\,\text{N}$.
}
\end{ex}

% ===== Câu 137 =====
% ID: AIQ_L10C3_FULL_0785
% Mức: NB
\begin{ex}
% ID: AIQ_L10C3_FULL_0785
% Mức: NB
Vật $3\,\text{kg}$ nằm cân bằng trên mặt phẳng nghiêng nhẵn góc 30°, được giữ bằng dây song song mặt phẳng. Lấy g= $10\,\text{m}$ /s^2. Tính lực căng dây.
\choice
{15 $N$}
{16 $N$}
{\True 15 $N$}
{30 $N$}
\loigiai{
Chọn C. T=mg sin30°=3·10·0,5= $15\,\text{N}$.
}
\end{ex}

% ===== Câu 138 =====
% ID: AIQ_L10C3_FULL_0786
% Mức: NB
\begin{ex}
% ID: AIQ_L10C3_FULL_0786
% Mức: NB
Biển quảng cáo khối lượng $4\,\text{kg}$ được treo cân bằng bằng hai dây giống nhau, đối xứng; mỗi dây hợp với phương ngang góc 30°. Lấy g= $10\,\text{m}$ /s^2. Tính lực căng mỗi dây.
\choice
{40 $N$}
{41 $N$}
{39 $N$}
{\True 40 $N$}
\loigiai{
Chọn D. 2 $T$ sin30°=mg nên T=mg= $40\,\text{N}$.
}
\end{ex}

% ===== Câu 139 =====
% ID: AIQ_L10C3_FULL_0787
% Mức: NB
\begin{ex}
% ID: AIQ_L10C3_FULL_0787
% Mức: NB
Vật $5\,\text{kg}$ nằm cân bằng trên mặt phẳng nghiêng nhẵn góc 30°, được giữ bằng dây song song mặt phẳng. Lấy g= $10\,\text{m}$ /s^2. Tính lực căng dây.
\choice
{\True 25 $N$}
{26 $N$}
{24 $N$}
{50 $N$}
\loigiai{
Chọn A. T=mg sin30°=5·10·0,5= $25\,\text{N}$.
}
\end{ex}

% ===== Câu 140 =====
% ID: AIQ_L10C3_FULL_0788
% Mức: TH
\begin{ex}
% ID: AIQ_L10C3_FULL_0788
% Mức: TH
Biển quảng cáo khối lượng $6\,\text{kg}$ được treo cân bằng bằng hai dây giống nhau, đối xứng; mỗi dây hợp với phương ngang góc 30°. Lấy g= $10\,\text{m}$ /s^2. Tính lực căng mỗi dây.
\choice
{60 $N$}
{\True 60 $N$}
{59 $N$}
{120 $N$}
\loigiai{
Chọn B. 2 $T$ sin30°=mg nên T=mg= $60\,\text{N}$.
}
\end{ex}

% ===== Câu 141 =====
% ID: AIQ_L10C3_FULL_0789
% Mức: TH
\begin{ex}
% ID: AIQ_L10C3_FULL_0789
% Mức: TH
Vật $2\,\text{kg}$ nằm cân bằng trên mặt phẳng nghiêng nhẵn góc 30°, được giữ bằng dây song song mặt phẳng. Lấy g= $10\,\text{m}$ /s^2. Tính lực căng dây.
\choice
{10 $N$}
{11 $N$}
{\True 10 $N$}
{20 $N$}
\loigiai{
Chọn C. T=mg sin30°=2·10·0,5= $10\,\text{N}$.
}
\end{ex}

% ===== Câu 142 =====
% ID: AIQ_L10C3_FULL_0790
% Mức: TH
\begin{ex}
% ID: AIQ_L10C3_FULL_0790
% Mức: TH
Biển quảng cáo khối lượng $3\,\text{kg}$ được treo cân bằng bằng hai dây giống nhau, đối xứng; mỗi dây hợp với phương ngang góc 30°. Lấy g= $10\,\text{m}$ /s^2. Tính lực căng mỗi dây.
\choice
{30 $N$}
{31 $N$}
{29 $N$}
{\True 30 $N$}
\loigiai{
Chọn D. 2 $T$ sin30°=mg nên T=mg= $30\,\text{N}$.
}
\end{ex}

% ===== Câu 143 =====
% ID: AIQ_L10C3_FULL_0791
% Mức: TH
\begin{ex}
% ID: AIQ_L10C3_FULL_0791
% Mức: TH
Vật $4\,\text{kg}$ nằm cân bằng trên mặt phẳng nghiêng nhẵn góc 30°, được giữ bằng dây song song mặt phẳng. Lấy g= $10\,\text{m}$ /s^2. Tính lực căng dây.
\choice
{\True 20 $N$}
{21 $N$}
{19 $N$}
{40 $N$}
\loigiai{
Chọn A. T=mg sin30°=4·10·0,5= $20\,\text{N}$.
}
\end{ex}

% ===== Câu 144 =====
% ID: AIQ_L10C3_FULL_0792
% Mức: VD
\begin{ex}
% ID: AIQ_L10C3_FULL_0792
% Mức: VD
Biển quảng cáo khối lượng $5\,\text{kg}$ được treo cân bằng bằng hai dây giống nhau, đối xứng; mỗi dây hợp với phương ngang góc 30°. Lấy g= $10\,\text{m}$ /s^2. Tính lực căng mỗi dây.
\choice
{50 $N$}
{\True 50 $N$}
{49 $N$}
{100 $N$}
\loigiai{
Chọn B. 2 $T$ sin30°=mg nên T=mg= $50\,\text{N}$.
}
\end{ex}

% ===== Câu 145 =====
% ID: AIQ_L10C3_FULL_0793
% Mức: VD
\begin{ex}
% ID: AIQ_L10C3_FULL_0793
% Mức: VD
Vật $6\,\text{kg}$ nằm cân bằng trên mặt phẳng nghiêng nhẵn góc 30°, được giữ bằng dây song song mặt phẳng. Lấy g= $10\,\text{m}$ /s^2. Tính lực căng dây.
\choice
{30 $N$}
{31 $N$}
{\True 30 $N$}
{60 $N$}
\loigiai{
Chọn C. T=mg sin30°=6·10·0,5= $30\,\text{N}$.
}
\end{ex}

% ===== Câu 146 =====
% ID: AIQ_L10C3_FULL_0794
% Mức: TH
\begin{ex}
% ID: AIQ_L10C3_FULL_0794
% Mức: TH
Vật $5\,\text{kg}$ nằm cân bằng trên mặt phẳng nghiêng nhẵn góc 30°, được giữ bằng dây song song mặt phẳng. Lấy g= $10\,\text{m}$ /s^2. Tính lực căng dây.
\choiceTF
{\True Giá trị cần tìm là $25\,\text{N}$.}
{Giá trị cần tìm là $26\,\text{N}$.}
{\True Phải xác định đúng các lực tác dụng và chọn chiều dương trước khi lập phương trình.}
{Có thể bỏ qua điều kiện áp dụng và chiều của lực mà kết quả vẫn luôn đúng.}
\loigiai{
\begin{itemchoice}
\item (Đúng) Giá trị cần tìm là $25\,\text{N}$. (Vì): T=mg sin30°=5·10·0,5= $25\,\text{N}$. b) Sai vì sai giá trị. c) Đúng. d) Sai vì phải kiểm tra điều kiện áp dụng và chiều lực
\item (Sai) Giá trị cần tìm là $26\,\text{N}$.
\item (Đúng) Phải xác định đúng các lực tác dụng và chọn chiều dương trước khi lập phương trình.
\item (Sai) Có thể bỏ qua điều kiện áp dụng và chiều của lực mà kết quả vẫn luôn đúng.
\end{itemchoice}
}
\end{ex}

% ===== Câu 147 =====
% ID: AIQ_L10C3_FULL_0795
% Mức: TH
\begin{ex}
% ID: AIQ_L10C3_FULL_0795
% Mức: TH
Biển quảng cáo khối lượng $6\,\text{kg}$ được treo cân bằng bằng hai dây giống nhau, đối xứng; mỗi dây hợp với phương ngang góc 30°. Lấy g= $10\,\text{m}$ /s^2. Tính lực căng mỗi dây.
\choiceTF
{\True Giá trị cần tìm là $60\,\text{N}$.}
{Giá trị cần tìm là $61\,\text{N}$.}
{\True Phải xác định đúng các lực tác dụng và chọn chiều dương trước khi lập phương trình.}
{Có thể bỏ qua điều kiện áp dụng và chiều của lực mà kết quả vẫn luôn đúng.}
\loigiai{
\begin{itemchoice}
\item (Đúng) Giá trị cần tìm là $60\,\text{N}$. (Vì): 2 $T$ sin30°=mg nên T=mg= $60\,\text{N}$. b) Sai vì sai giá trị. c) Đúng. d) Sai vì phải kiểm tra điều kiện áp dụng và chiều lực
\item (Sai) Giá trị cần tìm là $61\,\text{N}$.
\item (Đúng) Phải xác định đúng các lực tác dụng và chọn chiều dương trước khi lập phương trình.
\item (Sai) Có thể bỏ qua điều kiện áp dụng và chiều của lực mà kết quả vẫn luôn đúng.
\end{itemchoice}
}
\end{ex}

% ===== Câu 148 =====
% ID: AIQ_L10C3_FULL_0796
% Mức: VD
\begin{ex}
% ID: AIQ_L10C3_FULL_0796
% Mức: VD
Vật $2\,\text{kg}$ nằm cân bằng trên mặt phẳng nghiêng nhẵn góc 30°, được giữ bằng dây song song mặt phẳng. Lấy g= $10\,\text{m}$ /s^2. Tính lực căng dây.
\choiceTF
{\True Giá trị cần tìm là $10\,\text{N}$.}
{Giá trị cần tìm là $11\,\text{N}$.}
{\True Phải xác định đúng các lực tác dụng và chọn chiều dương trước khi lập phương trình.}
{Có thể bỏ qua điều kiện áp dụng và chiều của lực mà kết quả vẫn luôn đúng.}
\loigiai{
\begin{itemchoice}
\item (Đúng) Giá trị cần tìm là $10\,\text{N}$. (Vì): T=mg sin30°=2·10·0,5= $10\,\text{N}$. b) Sai vì sai giá trị. c) Đúng. d) Sai vì phải kiểm tra điều kiện áp dụng và chiều lực
\item (Sai) Giá trị cần tìm là $11\,\text{N}$.
\item (Đúng) Phải xác định đúng các lực tác dụng và chọn chiều dương trước khi lập phương trình.
\item (Sai) Có thể bỏ qua điều kiện áp dụng và chiều của lực mà kết quả vẫn luôn đúng.
\end{itemchoice}
}
\end{ex}

% ===== Câu 149 =====
% ID: AIQ_L10C3_FULL_0797
% Mức: VD
\begin{ex}
% ID: AIQ_L10C3_FULL_0797
% Mức: VD
Biển quảng cáo khối lượng $3\,\text{kg}$ được treo cân bằng bằng hai dây giống nhau, đối xứng; mỗi dây hợp với phương ngang góc 30°. Lấy g= $10\,\text{m}$ /s^2. Tính lực căng mỗi dây.
\choiceTF
{\True Giá trị cần tìm là $30\,\text{N}$.}
{Giá trị cần tìm là $31\,\text{N}$.}
{\True Phải xác định đúng các lực tác dụng và chọn chiều dương trước khi lập phương trình.}
{Có thể bỏ qua điều kiện áp dụng và chiều của lực mà kết quả vẫn luôn đúng.}
\loigiai{
\begin{itemchoice}
\item (Đúng) Giá trị cần tìm là $30\,\text{N}$. (Vì): 2 $T$ sin30°=mg nên T=mg= $30\,\text{N}$. b) Sai vì sai giá trị. c) Đúng. d) Sai vì phải kiểm tra điều kiện áp dụng và chiều lực
\item (Sai) Giá trị cần tìm là $31\,\text{N}$.
\item (Đúng) Phải xác định đúng các lực tác dụng và chọn chiều dương trước khi lập phương trình.
\item (Sai) Có thể bỏ qua điều kiện áp dụng và chiều của lực mà kết quả vẫn luôn đúng.
\end{itemchoice}
}
\end{ex}

% ===== Câu 150 =====
% ID: AIQ_L10C3_FULL_0798
% Mức: VD
\begin{ex}
% ID: AIQ_L10C3_FULL_0798
% Mức: VD
Vật $4\,\text{kg}$ nằm cân bằng trên mặt phẳng nghiêng nhẵn góc 30°, được giữ bằng dây song song mặt phẳng. Lấy g= $10\,\text{m}$ /s^2. Tính lực căng dây.
\choiceTF
{\True Giá trị cần tìm là $20\,\text{N}$.}
{Giá trị cần tìm là $21\,\text{N}$.}
{\True Phải xác định đúng các lực tác dụng và chọn chiều dương trước khi lập phương trình.}
{Có thể bỏ qua điều kiện áp dụng và chiều của lực mà kết quả vẫn luôn đúng.}
\loigiai{
\begin{itemchoice}
\item (Đúng) Giá trị cần tìm là $20\,\text{N}$. (Vì): T=mg sin30°=4·10·0,5= $20\,\text{N}$. b) Sai vì sai giá trị. c) Đúng. d) Sai vì phải kiểm tra điều kiện áp dụng và chiều lực
\item (Sai) Giá trị cần tìm là $21\,\text{N}$.
\item (Đúng) Phải xác định đúng các lực tác dụng và chọn chiều dương trước khi lập phương trình.
\item (Sai) Có thể bỏ qua điều kiện áp dụng và chiều của lực mà kết quả vẫn luôn đúng.
\end{itemchoice}
}
\end{ex}

% ===== Câu 151 =====
% ID: AIQ_L10C3_FULL_0799
% Mức: TH
\begin{ex}
% ID: AIQ_L10C3_FULL_0799
% Mức: TH
Biển quảng cáo khối lượng $3\,\text{kg}$ được treo cân bằng bằng hai dây giống nhau, đối xứng; mỗi dây hợp với phương ngang góc 30°. Lấy g= $10\,\text{m}$ /s^2. Tính lực căng mỗi dây.
\shortans{30}
\loigiai{
2 $T$ sin30°=mg nên T=mg= $30\,\text{N}$.
}
\end{ex}

% ===== Câu 152 =====
% ID: AIQ_L10C3_FULL_0800
% Mức: TH
\begin{ex}
% ID: AIQ_L10C3_FULL_0800
% Mức: TH
Vật $4\,\text{kg}$ nằm cân bằng trên mặt phẳng nghiêng nhẵn góc 30°, được giữ bằng dây song song mặt phẳng. Lấy g= $10\,\text{m}$ /s^2. Tính lực căng dây.
\shortans{20}
\loigiai{
$T=mg sin30°=4·10·0,5=20 N.$
}
\end{ex}

% ===== Câu 153 =====
% ID: AIQ_L10C3_FULL_0801
% Mức: VD
\begin{ex}
% ID: AIQ_L10C3_FULL_0801
% Mức: VD
Biển quảng cáo khối lượng $5\,\text{kg}$ được treo cân bằng bằng hai dây giống nhau, đối xứng; mỗi dây hợp với phương ngang góc 30°. Lấy g= $10\,\text{m}$ /s^2. Tính lực căng mỗi dây.
\shortans{50}
\loigiai{
2 $T$ sin30°=mg nên T=mg= $50\,\text{N}$.
}
\end{ex}

% ===== Câu 154 =====
% ID: AIQ_L10C3_FULL_0802
% Mức: VD
\begin{ex}
% ID: AIQ_L10C3_FULL_0802
% Mức: VD
Vật $6\,\text{kg}$ nằm cân bằng trên mặt phẳng nghiêng nhẵn góc 30°, được giữ bằng dây song song mặt phẳng. Lấy g= $10\,\text{m}$ /s^2. Tính lực căng dây.
\shortans{30}
\loigiai{
$T=mg sin30°=6·10·0,5=30 N.$
}
\end{ex}

% ===== Câu 155 =====
% ID: AIQ_L10C3_FULL_0803
% Mức: VD
\begin{ex}
% ID: AIQ_L10C3_FULL_0803
% Mức: VD
Biển quảng cáo khối lượng $2\,\text{kg}$ được treo cân bằng bằng hai dây giống nhau, đối xứng; mỗi dây hợp với phương ngang góc 30°. Lấy g= $10\,\text{m}$ /s^2. Tính lực căng mỗi dây.
\shortans{20}
\loigiai{
2 $T$ sin30°=mg nên T=mg= $20\,\text{N}$.
}
\end{ex}

% ===== Câu 156 =====
% ID: AIQ_L10C3_FULL_0804
% Mức: VDC
\begin{ex}
% ID: AIQ_L10C3_FULL_0804
% Mức: VDC
Vật $3\,\text{kg}$ nằm cân bằng trên mặt phẳng nghiêng nhẵn góc 30°, được giữ bằng dây song song mặt phẳng. Lấy g= $10\,\text{m}$ /s^2. Tính lực căng dây.
\shortans{15}
\loigiai{
$T=mg sin30°=3·10·0,5=15 N.$
}
\end{ex}

% ===== Câu 157 =====
% ID: AIQ_L10C3_FULL_0805
% Mức: VD
\begin{bt}
% ID: AIQ_L10C3_FULL_0805
% Mức: VD
Vật $6\,\text{kg}$ nằm cân bằng trên mặt phẳng nghiêng nhẵn góc 30°, được giữ bằng dây song song mặt phẳng. Lấy g= $10\,\text{m}$ /s^2. Tính lực căng dây.
\loigiai{
Phân tích lực, chọn hệ quy chiếu và áp dụng định luật hoặc quy tắc phù hợp. T=mg sin30°=6·10·0,5= $30\,\text{N}$.
}
\end{bt}

% ===== Câu 158 =====
% ID: AIQ_L10C3_FULL_0806
% Mức: VD
\begin{bt}
% ID: AIQ_L10C3_FULL_0806
% Mức: VD
Biển quảng cáo khối lượng $2\,\text{kg}$ được treo cân bằng bằng hai dây giống nhau, đối xứng; mỗi dây hợp với phương ngang góc 30°. Lấy g= $10\,\text{m}$ /s^2. Tính lực căng mỗi dây.
\loigiai{
Phân tích lực, chọn hệ quy chiếu và áp dụng định luật hoặc quy tắc phù hợp. 2 $T$ sin30°=mg nên T=mg= $20\,\text{N}$.
}
\end{bt}

% ===== Câu 159 =====
% ID: AIQ_L10C3_FULL_0807
% Mức: VDC
\begin{bt}
% ID: AIQ_L10C3_FULL_0807
% Mức: VDC
Vật $3\,\text{kg}$ nằm cân bằng trên mặt phẳng nghiêng nhẵn góc 30°, được giữ bằng dây song song mặt phẳng. Lấy g= $10\,\text{m}$ /s^2. Tính lực căng dây.
\loigiai{
Phân tích lực, chọn hệ quy chiếu và áp dụng định luật hoặc quy tắc phù hợp. T=mg sin30°=3·10·0,5= $15\,\text{N}$.
}
\end{bt}

% ===== Câu 160 =====
% ID: AIQ_L10C3_FULL_0808
% Mức: VDC
\begin{bt}
% ID: AIQ_L10C3_FULL_0808
% Mức: VDC
Biển quảng cáo khối lượng $4\,\text{kg}$ được treo cân bằng bằng hai dây giống nhau, đối xứng; mỗi dây hợp với phương ngang góc 30°. Lấy g= $10\,\text{m}$ /s^2. Tính lực căng mỗi dây.
\loigiai{
Phân tích lực, chọn hệ quy chiếu và áp dụng định luật hoặc quy tắc phù hợp. 2 $T$ sin30°=mg nên T=mg= $40\,\text{N}$.
}
\end{bt}

% ===== Câu 161 =====
% ID: AIQ_L10C3_FULL_0809
% Mức: VDC
\begin{bt}
% ID: AIQ_L10C3_FULL_0809
% Mức: VDC
Vật $5\,\text{kg}$ nằm cân bằng trên mặt phẳng nghiêng nhẵn góc 30°, được giữ bằng dây song song mặt phẳng. Lấy g= $10\,\text{m}$ /s^2. Tính lực căng dây.
\loigiai{
Phân tích lực, chọn hệ quy chiếu và áp dụng định luật hoặc quy tắc phù hợp. T=mg sin30°=5·10·0,5= $25\,\text{N}$.
}
\end{bt}

% ===== Câu 162 =====
% ID: AIQ_L10C3_FULL_0810
% Mức: VDC
\begin{bt}
% ID: AIQ_L10C3_FULL_0810
% Mức: VDC
Biển quảng cáo khối lượng $6\,\text{kg}$ được treo cân bằng bằng hai dây giống nhau, đối xứng; mỗi dây hợp với phương ngang góc 30°. Lấy g= $10\,\text{m}$ /s^2. Tính lực căng mỗi dây.
\loigiai{
Phân tích lực, chọn hệ quy chiếu và áp dụng định luật hoặc quy tắc phù hợp. 2 $T$ sin30°=mg nên T=mg= $60\,\text{N}$.
}
\end{bt}
